\subsection{Dimension and secant sequences}

Let A be a Noetherian ring, M a finitely generated A-module, S a subset of the radical of A, S the ideal of A generated by S, and a the annihilator of M. We have

\[
\dim_ {\mathrm{A}} (\mathrm{M}) = \dim (\mathrm{A} / \mathrm{a}) \leqslant \dim (\mathrm{A});\tag{6}
\]

Thus, if A is local, we have \(\dim_{\mathbf{A}}(\mathbf{M}) < +\infty\). Moreover, the support of the A-module M/SM is equal to V(a + S) by the corollary of Prop. 18 of II, § 4, no. 4, whence

\[
\dim_ {A} (M / S M) = \dim (A / (a + \mathfrak {S}))  .\tag{7}
\]

When S is finite, or when M is not reduced to 0, we have the inequality

\[
\dim_ {\mathbf {A}} (\mathbf {M} / \mathrm{SM}) \leqslant \dim_ {\mathbf {A}} (\mathbf {M}) \leqslant \operatorname{Card} (\mathbf {S}) + \dim_ {\mathbf {A}} (\mathbf {M} / \mathrm{SM});\tag{8}
\]

When S is finite, this follows from prop. 2 of no. 1 and from formulas (6) and (7) above, and the case where S is infinite is trivial.

DEFINITION 1. --- Let A be a Noetherian local ring, M a nonzero finitely generated A-module, and S a subset of the maximal ideal \(\mathfrak{m}_{\mathrm{A}}\) of A. One says that S is secant for M if one has

\[
\dim_ {A} (M) = \operatorname{Card} (S) + \dim_ {A} (M / S M).\tag{9}
\]

If S is secant for M, one has \(\operatorname{Card}(S) \leqslant \dim_{\mathbf{A}}(\mathbf{M})\) , so S is finite. One says that a family \((x_i)_{i \in I}\) of elements of \(\mathfrak{m}_{\mathbf{A}}\) is secant for M if one has

\[
\dim_ {\mathrm{A}} (\mathrm{M}) = \operatorname{Card} (\mathrm{I}) + \dim_ {\mathrm{A}} \left(\mathrm{M} / \sum_ {i \in \mathrm{I}} x _ {i} \mathrm{M}\right),\tag{10}
\]

that is, if \(i\mapsto x_{i}\) is a bijection from I onto a subset of \(\mathfrak{m}_{\mathbf{A}}\) secant for M.

We say that an element \(x\) of \(m_A\) is secant for \(M\) if \(\{x\}\) is a secant subset for \(M\) , that is, if one has

\[
\dim_ {\mathbf {A}} (\mathbf {M} / x \mathbf {M}) = \dim_ {\mathbf {A}} (\mathbf {M}) - 1.
\]

Remarks. --- 1) It follows from formulas (6) and (7) that S is secant for M if and only if it is secant for A/a, where a is the annihilator of M.

\begin{enumerate}
\def\labelenumi{\arabic{enumi})}
\setcounter{enumi}{1}
\tightlist
\item
  Let S and S′ be two disjoint subsets of \(m_{A}\) . For \(S \cup S'\) to be secant for M, it is necessary and sufficient that S be secant for M and S' be secant for \(M' = M/SM\). This follows from inequality (8) and the formula
\end{enumerate}

\[
\begin{array}{r l} \operatorname{Card} (S \cup S ^ {\prime}) + \dim_ {A} (M / (S M + S ^ {\prime} M)) - \dim_ {A} (M) = & \\ = (\operatorname{Card} (S) + \dim_ {A} (M / S M) - \dim_ {A} (M)) + & \\ + (\operatorname{Card} (S ^ {\prime}) + \dim_ {A} (M ^ {\prime} / S ^ {\prime} M ^ {\prime}) - \dim_ {A} (M ^ {\prime})) \end{array}
\]

\begin{enumerate}
\def\labelenumi{\arabic{enumi})}
\setcounter{enumi}{2}
\tightlist
\item
  Let \(x \in \mathfrak{m}_{\mathbf{A}}\) and \(n \geqslant 1\) be an integer. It is immediate that the modules \(\mathbf{M} / x\mathbf{M}\) and \(\mathbf{M} / x^n\mathbf{M}\) have the same support, hence the same dimension. Consequently, \(x\) is secant for \(\mathbf{M}\) if and only if \(x^n\) is secant for \(\mathbf{M}\). From this and Remark 2, we immediately deduce the following result: let \(x_1, ..., x_r\) be elements of \(\mathfrak{m}_{\mathbf{A}}\) and let \(n_1, ..., n_r\) be integers \(> 0\); then the sequence \((x_1, ..., x_r)\) is secant for \(\mathbf{M}\) if and only if the sequence \((x_1^{n_1}, ..., x_r^{n_r})\) is secant for \(\mathbf{M}\) .
\end{enumerate}

PROPOSITION 3. --- Let A be a Noetherian local ring and M a nonzero finitely generated A-module. For an element x of \(\mathfrak{m}_{\mathrm{A}}\) to be secant for M, it is necessary and sufficient that it belong to none of the minimal elements \(\mathfrak{p}\) of \(\operatorname{Supp}(M)\) such that \(\dim(\mathrm{A}/\mathfrak{p}) = \dim_{\mathrm{A}}(\mathrm{M})\), and it is sufficient that the homothety \(x_{\mathrm{M}}\) of ratio x on M be injective.

Let a be the annihilator of M. To say that x is secant for M means that x is secant for A/a, and the support of M consists of the prime ideals p of A such that \(a \subset p\); if \(x_{M}\) is injective, the image of x in A/a is not a zero divisor in A/a. Prop. 3 then follows from Cor. 2 of Prop. 2 of No. 1 applied to the ring A/a.

COROLLARY. --- Every sequence of elements of \(m_{A}\) which is completely secant for M (A, X, p. 157, def. 2) is secant for M.

Let \((x_{1},\ldots ,x_{r})\) be a sequence of elements of \(\mathfrak{m}_{\mathbf{A}}\) which is completely secant for M. Set \(\mathbf{M}_0 = \mathbf{M}\) and by induction \(\mathbf{M}_i = \mathbf{M}_{i - 1} / x_i\mathbf{M}_{i - 1}\) for \(1\leqslant i\leqslant r\) . By cor. 1 of A, X, p. 160, the homothety with ratio \(x_{i}\) in \(\mathbf{M}_{i - 1}\) is injective, whence \(\dim_{\mathbf{A}}(\mathbf{M}_i) = \dim_{\mathbf{A}}(\mathbf{M}_{i - 1}) - 1\) (for \(1\leqslant i\leqslant r\) ) (prop. 3). We therefore have

\[
\dim_ {\mathbf {A}} (\mathbf {M}) = r + \dim_ {\mathbf {A}} \left(\mathrm{M} / (x _ {1} \mathrm{M} + \dots + x _ {r} \mathrm{M})\right),
\]

hence the sequence \((x_{1},...,x_{r})\) is secant for M.

Remark 4. --- It is not true in general that a secant sequence for M is completely secant for M (p. 87, exerc. 6). We will study later the modules over a Noetherian local ring for which every secant sequence is completely secant.

THEOREM 1. --- Let A be a Noetherian local ring, M a nonzero finitely generated A-module, and S a subset of the maximal ideal \(\mathfrak{m}_{\mathbf{A}}\) of A.

\begin{enumerate}
\def\labelenumi{\alph{enumi})}
\item
  If M/SM is of finite length, we have Card(S) \(\geqslant\) dim \(_{A}\) (M); if S is a secant for M, then Card(S) \(\leqslant\) dim \(_{A}\) (M).
\item
  Every secant subset for M is contained in a maximal secant subset for M.
\item
  The following properties are equivalent:
\end{enumerate}

\begin{enumerate}
\def\labelenumi{(\roman{enumi})}
\item
  S is a maximal secant subset for M;
\item
  S is a secant subset for M and \(\operatorname{Card}(\mathbf{S}) = \dim_{\mathbf{A}}(\mathbf{M})\) ;
\item
  M/SM has finite length and Card(S) = dim\_A(M);
\item
  S is a secant subset for M and M/SM is of finite length.
\end{enumerate}

Since we have \(S \subset m_{A}\) , Nakayama's lemma shows that we have M/SM ≠ \{0\}, whence dim \(_{A}\) (M/SM) ≥ 0 with equality if and only if M/SM has finite length. Assertion \(a\) ) then follows from formulas (8) and (9), as well as the equivalence of properties (ii), (iii), and (iv).

The assertion \(b\) ) follows from the fact that the cardinal of every subset of \(\mathfrak{m}_{\mathbf{A}}\) secant for \(\mathbf{M}\) is bounded above by the integer \(\dim_{\mathbf{A}}(\mathbf{M})\) .

By \(a\) ), every subset secant for M, of cardinal equal to \(\dim_{\mathbf{A}}(\mathbf{M})\) , is maximal. It remains to prove that, if S is secant for M and if \(\operatorname{Card}(\mathbf{S}) < \dim_{\mathbf{A}}(\mathbf{M})\) , then S is not maximal. Let a be the annihilator of M, and B the local noetherian ring A/(a + SA). By cor. 2 of prop. 2 of n° 1, there exists an element \(x\) of \(\mathfrak{m}_{A}\) such that \(\dim(\mathrm{B}/x\mathrm{B}) = \dim(\mathrm{B}) - 1\); whence \(x \notin \mathrm{S}\) . By remark 2, the subset S ∪ \(\{x\}\) of \(\mathfrak{m}_{A}\) is secant for A/a, hence for M by remark 1.

COROLLARY. --- The dimension of M is the smallest of the integers \(d \geqslant 0\) for which there exists a sequence \((x_{1}, ..., x_{d})\) of elements of \(m_{A}\) such that the A-module M/ \(\sum_{i=1}^{d} x_{i}M\) has finite length.

Since \(\varnothing\) is a secant subset for M, Theorem 1, \(b\) ) proves the existence of a maximal secant sequence for M, namely \((x_{1},\ldots,x_{d})\) . But then we have \(d=\dim_{\mathbf{A}}(\mathbf{M})\) and the A-module \(\mathbf{M}/\sum_{i=1}^{d}x_{i}\mathbf{M}\) has finite length by property (iii) of Thm. 1, \(c\) ). Conversely, if \((x_{1}',\ldots,x_{d'}')\) is a sequence of elements of \(\mathfrak{m}_{\mathbf{A}}\) such that the A-module \(\mathbf{M}/\sum_{j=1}^{d'}x_{j}'\mathbf{M}\) is of finite length, we have \(d'\geqslant\dim_{\mathbf{A}}(\mathbf{M})\) by th. 1, \(a\) ).

Recall (III, § 3, n° 2, def. 1) that an ideal q of a Noetherian local ring A is an ideal of definition of A if the q-adic and m\_A-adic topologies of A coincide.

Lemma 2. --- Let A be a Noetherian local ring and q an ideal of A. The following conditions are equivalent:

\begin{enumerate}
\def\labelenumi{(\roman{enumi})}
\item
  q is an ideal of definition of A;
\item
  there exists an integer \(n \geqslant 0\) such that \(\mathfrak{m}_{\mathbf{A}}^{n} \subset \mathfrak{q} \subset \mathfrak{m}_{\mathbf{A}}\) ;
\item
  we have q ≠ A and A/q is an A-module of finite length;
\item
  V(q) is equal to \(\{m_{A}\}\) (in other words, \(m_{A}\) is the only prime ideal of A containing q).
\end{enumerate}

Indeed, the equivalence of (i), (ii), and (iv) was proved in III, § 2, no. 5, and the equivalence of (i) and (iii) follows from IV, § 2, no. 5, cor. 2 to prop. 9.

The corollary of Thm. 1 therefore allows us to state the following scholium:

SCHOLIUM. --- The dimension of a Noetherian local ring A is the smallest integer \(d \geqslant 0\) for which there exists an ideal of definition of A generated by d elements.

\subsection{First applications}

Let A be a Noetherian ring, V = Spec(A) its spectrum. In this number, we call a hypersurface in V any subset of the form \footnote{Recall that \(V(x)\) consists of the prime ideals of \(A\) containing \(x\).} V(x) with \(x \in A\) .

PROPOSITION 4. --- Let X be a closed subset of V, and let \(H_{1}, \ldots, H_{m}\) be hypersurfaces in V. Set \(X' = X \cap H_{1} \cap \ldots \cap H_{m}\) .

\begin{enumerate}
\def\labelenumi{\alph{enumi})}
\tightlist
\item
  For every closed subset Y of V contained in X′, we have
\end{enumerate}

\[
\operatorname{codim} (\mathrm{Y}, \mathrm{X} ^ {\prime}) \geqslant \operatorname{codim} (\mathrm{Y}, \mathrm{X}) - m.
\]

\begin{enumerate}
\def\labelenumi{\alph{enumi})}
\setcounter{enumi}{1}
\item
  We have codim(Z, X) ≤ m for every irreducible component Z of X'. If X' is nonempty, then codim(X', X) ≤ m.
\item
  If Z is a closed irreducible subset of V contained in X such that \(\text{codim}(Z, X) \leqslant m\), there exist hypersurfaces \(H_{1}^{\prime}, \ldots, H_{m}^{\prime}\) such that Z is an irreducible component of \(X \cap H_{1}^{\prime} \cap \ldots \cap H_{m}^{\prime}\) .
\end{enumerate}

Let \(\mathfrak{a}\) be an ideal of A and \(x_{1},\ldots,x_{m}\) elements of A such that \(\mathrm{X}=\mathrm{V}(\mathfrak{a})\) and \(\mathrm{H}_{i}=\mathrm{V}(x_{i})\) for \(1\leqslant i\leqslant m\) . Let Z be a closed irreducible subset of V contained in X; there exists a prime ideal p of A containing \(\mathfrak{a}\) and such that \(\mathrm{Z}=\mathrm{V}(\mathfrak{p})\) .

Suppose first that Z is contained in X' and denote by \(\xi_{i}\) the image of \(x_{i}\) in the Noetherian local ring \(\mathbf{B} = \mathbf{A}_{\mathfrak{p}} / \mathfrak{a}\mathbf{A}_{\mathfrak{p}}\) . By prop. 7, b) of § 1, no. 3, we have

\[
\operatorname{codim} (Z, X) = \dim (B), \quad \operatorname{codim} (Z, X ^ {\prime}) = \dim (B / (\xi_ {1} B + \dots + \xi_ {m} B)).
\]

By prop. 2 of no. 1, we therefore have

\[
\operatorname{codim} (Z, X ^ {\prime}) \geqslant \operatorname{codim} (Z, X) - m.\tag{11}
\]

If Z is an irreducible component of X', we have codim(Z, X') = 0, whence codim(Z, X) ≤ m; this proves b). We prove a) by taking, in the two members of (11), the lower bound over the set of irreducible components Z of Y.

Conversely, suppose we have \(\operatorname{codim}(\mathbf{Z},\mathbf{X})\leqslant m\), that is, \(\dim (\mathbf{B})\leqslant m\). Since every element of \(\mathbf{A}_{\mathfrak{p}}\) is the product of an invertible element of \(\mathbf{A}_{\mathfrak{p}}\) by the image of an element of A, the scholium of No. 2 proves the existence of elements \(x_1^{\prime},\ldots ,x_m^\prime\) of A whose images in B generate an ideal of definition of B. Set \(\mathrm{H}_i^{\prime} = \mathrm{V}(x_i^{\prime})\) for \(1\leqslant i\leqslant m\). It is clear that Z is an irreducible component of \(\mathbf{X}\cap \mathbf{H}_1^{\prime}\cap \dots \cap \mathbf{H}_m^{\prime}\) .

COROLLARY 1. --- Let H be a nonempty hypersurface in V. The codimension of H in V is equal to 0 or 1.

We have codim(H, V) = 1 if and only if H contains no irreducible component of V. If this is the case, then all irreducible components of H have codimension 1 in V.

For every irreducible component Z of H, we have

\[
0 \leqslant \operatorname{codim} (Z, V) \leqslant 1
\]

by prop. 4, b), and one has codim(Z, V) = 0 if and only if Z is an irreducible component of V. By definition,

\[
\operatorname{codim} (\mathrm{H}, \mathrm{V}) = \inf _ {\mathrm{Z}} \operatorname{codim} (\mathrm{Z}, \mathrm{V})
\]

where Z runs through the set of irreducible components of H. Cor. 1 follows immediately from these remarks.

COROLLARY 2. --- Let X be a closed irreducible subset of V and H a hypersurface in V. Only three cases are possible:

\begin{enumerate}
\def\labelenumi{\arabic{enumi})}
\item
  we have \(\mathbf{X}\subset \mathbf{H}\)
\item
  the set \(\mathbf{X} \cap \mathbf{H}\) is nonempty and each of its irreducible components \(Z\) satisfies \(\operatorname{codim}(\mathbf{Z}, \mathbf{X}) = 1\) ;
\item
  The set X ∩ H is empty.
\end{enumerate}

Suppose \(X' = X \cap H\) is nonempty and distinct from \(X\); every irreducible component \(Z\) of \(X'\) is distinct from \(X\) and satisfies codim \((Z, X) \leqslant 1\) by prop. 4, b). Cor. 2 follows immediately from this.

COROLLARY 3. --- If A is factorial (VII, § 3, no. 3, prop. 2), the prime ideals of height 1 of A are the principal ideals generated by the extremal elements of A. If moreover A is local, then dim(A/p) = dim(A) - 1 for every prime ideal p of height 1 of A.

Let \(x\) be an extremal element of A. Then \(Ax\) is a prime ideal because \(x\) is extremal, of height 1 because A is integral (no. 1, prop. 1). Let \(p\) be a prime ideal of height 1 of A. Then \(V(p)\) is an irreducible component of a hypersurface \(V(Ax)\) for a suitable \(x\) (prop. 4, c)). Let \(x = \prod_{i} y_{i}^{n_i}\) be a decomposition of \(x\) into products of extremal

elements such that \(y_{i}\) and \(y_{j}\) are coprime if \(i \neq j\). The irreducible components of \(V(Ax)\) are the \(V(Ay_{i})\). Therefore \(p = Ay_{i}\) for a suitable i. The last assertion follows from cor. 2 to prop. 2 of no. 1.

Remarks. --- 1) Suppose that A is a local, Noetherian, integral domain of dimension d; let x be a nonzero element of \(m_{A}\) and let \(H = V(x)\) . By cor. 2 of prop. 4, every irreducible component of H has codimension 1 in X, hence dimension \(\leqslant d - 1\) (§ 1, no. 2, prop. 3). By cor. 2 of prop. 2 of no. 1, H has dimension d - 1, and one of these components therefore has dimension d - 1; all of them do if A is catenary. However, in general there may exist an irreducible component of H of dimension \textless{} d - 1 (cf. p. 87, exercise 7).

\begin{enumerate}
\def\labelenumi{\arabic{enumi})}
\setcounter{enumi}{1}
\item
  Let \(x_1, \ldots, x_n\) be elements of A, and set \(\mathrm{H}_i = \mathrm{V}(x_i)\) for \(1 \leqslant i \leqslant n\). Suppose there exists a finitely generated A-module M with support \(\mathrm{V} = \operatorname{Spec}(\mathrm{A})\) such that \((x_1, \ldots, x_n)\) is a completely secant sequence for M. Then every irreducible component of \(\mathrm{H}_1 \cap \ldots \cap \mathrm{H}_n\) is of codimension \(n\) in V: this follows easily from the corollary of Prop. 3 of No. 2.
\item
  If the ideal a of the Noetherian ring A is generated by m elements, then ht(a) ≤ m. This follows immediately from prop. 4.
\end{enumerate}

PROPOSITION 5. --- Let A be a Noetherian ring and \(\mathfrak{p} \subset \mathfrak{q}\) a non-saturated chain of prime ideals of A. The set E of prime ideals \(\mathfrak{r}\) of A such that \(\mathfrak{p} \subset \mathfrak{r} \subset \mathfrak{q}\) is an infinite chain. We have \(\bigcup \mathfrak{r} = \mathfrak{q}\) and \(\bigcap \mathfrak{r} = \mathfrak{p}\) .

By replacing A with A/p, we reduce to the case where p = \{0\}.

By lemma 1 of Section 1, we have q = \(\bigcup_{r \in E} r\) and prop. 2 of II, § 1, no. 1 shows that E is infinite.

Let \(y \neq 0\) be an element of \(\bigcap_{r \in E} r\). The height of q is finite ( \(n^{\circ}1\) , cor. 1 of prop. 2), and one has \(ht(q) \geqslant 2\) by hypothesis. There therefore exists a prime ideal \(q' \subset q\) of height 2. The first part of the proof applied to \(q'\) shows that the set \(E'\) of prime ideals of height 1 contained in \(q'\) is infinite; each of these ideals contains y by hypothesis. Now the Noetherian local ring \(B = A_{q'} / yA_{q'}\) is of dimension 1 according to cor. 2 of prop. 2 of n° 1. For every \(r \in E'\) , the prime ideal \(r/yA_{q'}\) of B is therefore minimal; consequently, the Noetherian ring B has infinitely many minimal prime ideals, which is absurd (II, § 4, n° 3, cor. 3 of prop. 14). One therefore has \(\bigcap_{r \in E} r = \{0\}\) .

PROPOSITION 6.--- Let A be a Noetherian ring of dimension \(\geqslant 2\) , and \(h\) be an integer such that \(0 < h < \dim(A)\) .

\begin{enumerate}
\def\labelenumi{\alph{enumi})}
\item
  A possesses infinitely many prime ideals of height h.
\item
  If A is of finite dimension, it possesses infinitely many prime ideals p such that \(\operatorname{ht}(\mathfrak{p}) = h\) and \(\dim(\mathbf{A}/\mathfrak{p}) = \dim(\mathbf{A}) - h\) .
\end{enumerate}

Since the dimension of A is the upper bound of the heights of the prime ideals of A (§ 1, n° 3, prop. 8), since every prime ideal of A is of finite height (n° 1, cor. 1 of prop. 2) and since \(h < \dim(A)\) , there exists an integer \(n > h\) and a prime ideal \(p\) of height \(n\) , hence a chain \(p_0 \subset \ldots \subset p_n = p\) of prime ideals of length \(n\) . We have \(ht(p_i) = i\) for \(0 \leqslant i \leqslant n\) , whence \(ht(r) = h\) for every prime ideal \(r\) of A such that \(p_{h-1} \subset r \subset p_{h+1}\). The set E of these ideals is infinite according to prop. 5, whence \(a\) ).

If A is of finite dimension, one may suppose that one has \(n = \dim(A)\) in what precedes. For every prime ideal \(r \in E\) , one has \(\operatorname{ht}(r) = h\) , whence \(\dim(A/r) \leqslant n - h\) , and since \(r \subset p_{h+1} \subset \ldots \subset p_n\) is a chain of length \(n - h\) , one has \(\dim(A/r) = n - h\) .

There exist non-Noetherian integral domains of dimension 2 possessing only one prime ideal of height 1, for example a valuation ring of height 2 (VI, § 4, n° 4, prop. 5).

\subsection{Changes of rings}

PROPOSITION 7.--- Let \(\rho: A \to B\) be a local homomorphism of Noetherian local rings, M a finitely generated A-module and N a finitely generated B-module. Put \(\overline{B} = B \otimes_A \kappa_A = B / \rho(m_A).B\) and \(\overline{N} = N \otimes_B \overline{B} = N / \rho(m_A).N\) . We have

\[
\dim_ {B} (M \otimes_ {A} N) \leqslant \dim_ {A} (M) + \dim_ {B} (\overline {{N}}),\tag{12}
\]

and there is equality if N is flat over A.

One may suppose M and N non-zero. According to the corollary of th. 1 (n° 2), there exists a subset S (resp. T) of \(m_A\) (resp. \(m_B\) ) such that \(\text{Card}(S) = \dim_A(M)\) (resp. \(\text{Card}(T) = \dim_{\overline{B}}(\overline{N})\) ) and that M/SM is an A-module of finite length (resp. \(\overline{N}/\overline{T}\overline{N}\) is a B-module of finite length). One has \(\rho(S) \subset m_B\) . Let E be the B-module \(M \otimes_A N\) ; the B-module \(E/(\rho(S).E + T.E)\) is isomorphic to M/SM \(\otimes_A N/TN\) , hence is of finite length: according to prop. 18 and 19 of II, § 4, n° 4, one has

\[
\begin{array}{r l} \operatorname{Supp} (\mathrm{M} / \mathrm{SM} \otimes_ {\mathrm{A}} \mathrm{N} / \mathrm{TN}) & = \operatorname{Supp} (\mathrm{N} / \mathrm{TN}) \cap {} ^ {a} \rho^ {- 1} (\operatorname{Supp} (\mathrm{M} / \mathrm{SM})) \\ & = \operatorname{Supp} (\mathrm{N} / \mathrm{TN}) \cap {} ^ {a} \rho^ {- 1} (\operatorname{Supp} (\kappa_ {\mathrm{A}})) \\ & = \operatorname{Supp} (\mathrm{N} / \mathrm{TN} \otimes_ {\mathrm{A}} \kappa_ {\mathrm{A}}) = \operatorname{Supp} (\overline {{\mathrm{N}}} / \mathrm{T} \overline {{\mathrm{N}}}) = \left\{\mathfrak {m} _ {\mathrm{B}} \right\}. \end{array}
\]

It follows, according to th. 1, that one has the inequality

\[
\dim_ {\mathrm{B}} (\mathrm{E}) \leqslant \operatorname{Card} (\mathrm{S}) + \operatorname{Card} (\mathrm{T}) = \dim_ {\mathrm{A}} (\mathrm{M}) + \dim_ {\overline {{\mathrm{B}}}} (\overline {{\mathrm{N}}}).
\]

Suppose now N flat over A, and let us show that there is equality in formula (12). Let a (resp. b) be the annihilator of M (resp. N). According to prop. 18 and 19 of II, § 4, n° 4, one has

\[
\begin{array}{r l} \operatorname{Supp} (E) & = \operatorname{Supp} (N) \cap {} ^ {a} \rho^ {- 1} (\operatorname{Supp} (M)) \\ & = V (b) \cap {} ^ {a} \rho^ {- 1} (V (a)) = V (b + \rho (a). B). \end{array}
\]

One consequently has

\[
\dim_ {B} (M \otimes_ {A} N) = \dim (B / (b + \rho (a). B))\tag{13}
\]

and also

\[
\dim_ {\mathbf {A}} (\mathbf {M}) + \dim_ {\overline {{\mathbf {B}}}} (\overline {{\mathbf {N}}}) = \dim (\mathrm{A} / \mathfrak {a}) + \dim (\mathbf {B} / (\mathfrak {b} + \rho (m _ {\mathbf {A}}). \mathbf {B}))  .\tag{14}
\]

Put A' = A/a and B' = B/(b + ρ(a).B) and let ρ':A' → B' be the local homomorphism deduced from ρ by passing to quotients. Since the annihilator of N is b, N is a finitely generated B/b-module with support equal to Spec(B/b), and flat over A. The homomorphism A → B/b deduced from ρ therefore possesses property (PM) of § 2, n° 1 (loc. cit., remark 3); by extension of scalars (loc. cit., prop. 1, a)), one deduces that ρ' possesses property (PM). According to loc. cit., prop. 2, one therefore has

\[
\dim (\mathbf {B} ^ {\prime}) \geqslant \dim (\mathbf {A} ^ {\prime}) + \dim (\mathbf {B} ^ {\prime} / \rho^ {\prime} (m _ {\mathbf {A} ^ {\prime}}). \mathbf {B} ^ {\prime})\tag{15}
\]

and since the ring \(\mathbf{B}'/\rho'(\mathfrak{m}_{\mathbf{A}'})\mathbf{B}'\) is isomorphic to \(\mathbf{B}/(\mathfrak{b} + \rho(\mathfrak{m}_{\mathbf{A}})\mathbf{B})\), our assertion follows from formulas (13), (14) and (15).

COROLLARY 1. --- Let \(\rho: A \to B\) be a local homomorphism of Noetherian local rings.

\begin{enumerate}
\def\labelenumi{\alph{enumi})}
\tightlist
\item
  We have
\end{enumerate}

\[
\dim (\mathbf {B}) \leqslant \dim (\mathbf {A}) + \dim (\mathbf {B} \otimes_ {\mathbf {A}} \kappa_ {\mathbf {A}})\tag{16}
\]

with equality if ρ makes B into a flat A-module.

\begin{enumerate}
\def\labelenumi{\alph{enumi})}
\setcounter{enumi}{1}
\tightlist
\item
  Suppose that \(\rho\) makes B into a flat A-module, and let S be a subset of \(\mathfrak{m}_{\mathbf{A}}\) . Then S is secant for A if and only if \(\rho(\mathbf{S})\) is secant for B.
\end{enumerate}

The assertion \(a)\) is the particular case \(\mathbf{M} = \mathbf{A}\) , \(\mathbf{N} = \mathbf{B}\) from prop. 7.

Let us prove b). Under the stated hypotheses, we have

\[
\dim (\mathbf {B}) = \dim (\mathbf {A}) + \dim (\overline {{\mathbf {B}}})\tag{17}
\]

with \(\overline{\mathbf{B}} = \mathbf{B} \otimes_{\mathbf{A}} \kappa_{\mathbf{A}} = \mathrm{B} / \rho(\mathfrak{m}_{\mathbf{A}}) . \mathbf{B}\) . Since \(\rho\) is injective (I, § 3, no. 5, prop. 9), we have

\[
\operatorname{Card} (\rho (S)) = \operatorname{Card} (S).\tag{18}
\]

Finally, B′ = B/ρ(S). B is a flat module over A′ = A/SA, hence

\[
\dim (\mathrm{B} / \rho (\mathrm{S}). \mathrm{B}) = \dim (\mathrm{A} / \mathrm{SA}) + \dim (\overline {{\mathrm{B}}})\tag{19}
\]

since \(\mathbf{B}^{\prime}/\rho(\mathfrak{m}_{\mathbf{A}^{\prime}})\mathbf{B}^{\prime}\) is isomorphic to \(\overline{B}\). Our assertion follows immediately from formulas (17), (18), and (19).

COROLLARY 2. — Let ρ : A → B be a homomorphism of Noetherian rings. Then

\[
\dim (\mathbf {B}) \leqslant \dim (\mathbf {A}) + \sup _ {\mathfrak {p} \in \operatorname{Spec} (\mathbf {A})} \dim (\mathbf {B} \otimes_ {\mathbf {A}} \kappa (\mathfrak {p})).
\]

Let q be a prime ideal of B and \(\mathfrak{p} = \rho^{-1}(\mathfrak{q})\) . By cor. 1, we have \(\dim(B_{\mathfrak{q}}) \leqslant \dim(A_{\mathfrak{p}}) + \dim(B_{\mathfrak{q}} \otimes_{A} \kappa(\mathfrak{p}))\). From this we deduce (prop. 6, b) of § 1, no. 3) the inequality \(\dim(B_{\mathfrak{q}}) \leqslant \dim(A) + \dim(B \otimes_{A} \kappa(\mathfrak{p}))\) , and we conclude by prop. 8 of § 1, no. 3.

COROLLARY 3. --- Let A be a Noetherian ring and n an integer ≥ 0. Then

\[
\dim (\mathrm{A} [ \mathrm{X} _ {1},..., \mathrm{X} _ {n} ]) = \dim (\mathrm{A}) + n.
\]

Set \(\mathbf{B} = \mathbf{A}[\mathbf{X}_1, \ldots, \mathbf{X}_n]\). For every prime ideal \(\mathfrak{p}\) of \(\mathbf{A}\) , the ring \(\mathbf{B} \otimes_{\mathbf{A}} \kappa(\mathfrak{p})\) is a polynomial ring in \(n\) variables over a field, hence has dimension \(n\) ( \(\S 2\) , no 4, cor. 1 to th. 3). By cor. 2, we have \(\dim(\mathbf{B}) \leqslant \dim(\mathbf{A}) + n\); the reverse inequality follows from Example 4 of \(\S 1\) , no. 3.

COROLLARY 4. --- Let \(\rho: A \to B\) be a homomorphism of Noetherian rings, and a an ideal of A. We have the inequality

\[
\operatorname{ht} (\rho (\mathfrak {a}). \mathbf {B}) \leqslant \operatorname{ht} (\mathfrak {a})\tag{20}
\]

if the mapping \(^a\rho : \text{Spec}(\mathbf{B}) \to \text{Spec}(\mathbf{A})\) is surjective. If B is a faithfully flat A-module, then we have \(\operatorname{ht}(\rho(\mathfrak{a}).\mathbf{B}) = \operatorname{ht}(\mathfrak{a})\) .

If B is faithfully flat over A, then \(^{a}\rho\) is surjective (II, § 2, n° 5, cor. 4 to prop. 11) and one has ht(a) ≤ ht(ρ(a).B) (§ 2, n° 1, corollary of prop. 2).

It therefore remains to prove inequality (20) under the assumption that \(^a\rho\) is surjective. Let \(\mathfrak{p}\) be a prime ideal of A such that \(\mathfrak{a} \subset \mathfrak{p}\) and \(\mathrm{ht}(\mathfrak{a}) = \mathrm{ht}(\mathfrak{p})\) ( \(\S 1, n^o 3, \text{prop. } 7\) ). Let \(\overline{\mathbf{B}} = \mathbf{B} \otimes_{\mathbf{A}} \kappa(\mathfrak{p})\) and denote by \(h\) the canonical homomorphism from B to \(\overline{\mathbf{B}}\) . If \(X = ^a\rho^{-1}(\mathfrak{p})\) is the nonempty set of prime ideals of B lying over \(\mathfrak{p}\) , we know ( \(\S 2, n^o 1, \text{lemme 1}\) ) that the map \(^ah\) is a homeomorphism from \(\operatorname{Spec}(\overline{\mathbf{B}})\) onto the subspace X of \(\operatorname{Spec}(\mathbf{B})\) . Let \(\mathfrak{q}\) be the image under \(^ah\) of a minimal prime ideal of \(\overline{\mathbf{B}}\) ; one has \(\dim(\mathbf{B}_{\mathfrak{q}} \otimes_{\mathbf{A}} \kappa(\mathfrak{p})) = 0\) and Cor. 1 implies the inequality \(\dim(\mathbf{B}_{\mathfrak{q}}) \leqslant \dim(\mathbf{A}_{\mathfrak{p}})\) . Finally,

\[
\operatorname{ht} (\rho (a). B) \leqslant \operatorname{ht} (q) = \dim (B _ {q}) \leqslant \dim (A _ {p}) = \operatorname{ht} (p) = \operatorname{ht} (a),
\]

whence the corollary.

PROPOSITION 8.--- Let A be a Noetherian ring, a an ideal of A, M a finitely generated A-module, \(\hat{\mathbf{A}}\) and \(\hat{\mathbf{M}}\) the separated completions of A and M respectively for the a-adic topology.

\begin{enumerate}
\def\labelenumi{\alph{enumi})}
\item
  Let m be a prime ideal of A containing a and \(\hat{\mathfrak{m}} = \mathfrak{m}\hat{\mathbf{A}}\) . Then \(\hat{\mathfrak{m}}\) is a prime ideal of \(\hat{\mathbf{A}}\) and one has \(\dim_{\hat{\mathbf{A}}_{\hat{\mathfrak{m}}}}(\hat{\mathbf{M}}_{\hat{\mathfrak{m}}}) = \dim_{\mathbf{A}_{\mathfrak{m}}}(\mathbf{M}_{\mathfrak{m}})\) .
\item
  We have \(\dim_{\hat{\mathbf{A}}}(\hat{\mathbf{M}}) = \sup_{\mathfrak{m}} \dim_{\mathbf{A}_{\mathfrak{m}}}(\mathbf{M}_{\mathfrak{m}})\) , where \(\mathfrak{m}\) ranges over the set of prime ideals
\end{enumerate}

(resp. maximal ideals) of A containing a. In particular, one has \(\dim_{\hat{\mathbf{A}}}(\hat{\mathbf{M}}) \leqslant \dim_{\mathbf{A}}(\mathbf{M})\) .

\begin{enumerate}
\def\labelenumi{\alph{enumi})}
\item
  Since \(\hat{\mathbf{A}}/\hat{\mathfrak{m}}\) is identified with \(\mathbf{A}/\mathfrak{m}\) , \(\hat{\mathfrak{m}}\) is a prime ideal of \(\hat{\mathbf{A}}\) . By th. 3 of III, § 3, n° 4, \(\hat{\mathbf{A}}\) is flat over \(\mathbf{A}\) , hence \(\hat{\mathbf{A}}_{\hat{\mathfrak{m}}}\) is flat over \(\mathbf{A}_{\mathfrak{m}}\) . Moreover, the canonical map from \(\mathbf{A}\) into \(\hat{\mathbf{A}}\) induces an isomorphism of \(\mathbf{A}/\mathfrak{a}\) onto \(\hat{\mathbf{A}}/\mathfrak{a}\hat{\mathbf{A}}\) , hence also an isomorphism of \(\mathbf{A}_{\mathfrak{m}}/\mathfrak{mA}_{\mathfrak{m}}\) onto \(\hat{\mathbf{A}}_{\hat{\mathfrak{m}}}/\mathfrak{m}\hat{\mathbf{A}}_{\hat{\mathfrak{m}}}\) . We conclude by applying prop. 7 to the rings \(\mathbf{A}_{\mathfrak{m}}\) and \(\hat{\mathbf{A}}_{\hat{\mathfrak{m}}}\) and to the modules \(\mathbf{M}_{\mathfrak{m}}\) and \(\hat{\mathbf{A}}_{\hat{\mathfrak{m}}}\); the module \(\mathbf{M}_{\mathfrak{m}} \otimes_{\mathbf{A}_{\mathfrak{m}}} \hat{\mathbf{A}}_{\hat{\mathfrak{m}}}\) is isomorphic to \(\hat{\mathbf{M}}_{\hat{\mathfrak{m}}}\) (III, loc. cit. and prop. 8).
\item
  By prop. 8 of III, § 3, n° 4, the map m ↦ m̂ is a bijection from the set of maximal ideals of A containing a onto the set of maximal ideals of Â. Assertion b) follows from this and from prop. 9 of § 1, n° 4.
\end{enumerate}

COROLLARY 1. --- Let A be a Zariski ring (III, § 3, n° 3, def. 2). For every finitely generated A-module M, we have \(\dim_{\hat{\mathbf{A}}}(\hat{\mathbf{M}}) = \dim_{\mathbf{A}}(\mathbf{M})\) .

Indeed, the topology of A is the a-adic topology, where a is an ideal contained in the radical of A (loc. cit.), that is, contained in every maximal ideal m of A. It therefore suffices to apply assertion b) of prop. 8.

COROLLARY 2. --- Let A be a Noetherian ring, a an ideal of A, and \(\hat{A}\) the separated completion of A for the a-adic topology. We have dim( \(\hat{A}\) ) \(\leqslant\) dim(A), with equality when A is local and \(a\) is distinct from A.

COROLLARY 3. --- Let A be a Noetherian ring and n an integer ≥ 0. Then

\[
\dim (\mathrm{A} [ [ \mathrm{X} _ {1},..., \mathrm{X} _ {n} ] ]) = \dim (\mathrm{A}) + n.
\]

The ring A{[}{[}X₁, \ldots, Xₙ{]}{]} is the separated completion of the polynomial ring A{[}X₁, \ldots, Xₙ{]} for the \(\mathfrak{a}\)-adic topology, where \(\mathfrak{a}\) is the ideal generated by X₁, \ldots, Xₙ; we therefore have

\[
\dim (\mathrm{A} [ [ \mathrm{X} _ {1},..., \mathrm{X} _ {n} ] ]) \leqslant \dim (\mathrm{A} [ \mathrm{X} _ {1},..., \mathrm{X} _ {n} ])
\]

by cor. 2. We also have

\[
\dim (\mathrm{A} [ \mathrm{X} _ {1},..., \mathrm{X} _ {n} ]) = \dim (\mathrm{A}) + n
\]

by cor. 3 of prop. 7. Finally, we have

\[
\dim (\mathrm{A}) + n \leqslant \dim (\mathrm{A} [ [ \mathrm{X} _ {1},..., \mathrm{X} _ {n} ] ])
\]

by example 4 of § 1, n° 3.

Remarks. --- 1) Let A be a Noetherian ring and a an ideal of A. Suppose A is separated and complete for the a-adic topology, and consider the ring R = A \{ X \(_{1}\) , \ldots, X \(_{n}\) \} of restricted formal series (III, § 4, n° 2, def. 2). We have dim(R) = dim(A) + n. Indeed, R is the completion of the ring B = A{[}X \(_{1}\) , \ldots, X \(_{n}\) {]} for the aB-adic topology, whence dim(R) ≤ dim(A{[}X \(_{1}\) , \ldots, X \(_{n}\) {]}) = dim(A) + n. The reverse inequality is proved as in the case of formal series.

\begin{enumerate}
\def\labelenumi{\arabic{enumi})}
\setcounter{enumi}{1}
\tightlist
\item
  Let A be a Noetherian local ring, identified with a subring of its completion \(\hat{\mathbf{A}}\) , and let B be a subring of \(\hat{\mathbf{A}}\) containing A. Suppose that B is Noetherian local and that we have \(m_{\mathbf{A}}B = m_{\mathbf{B}}\) . Then the canonical injection of A into B extends to an isomorphism of \(\hat{\mathbf{A}}\) onto the completion \(\hat{\mathbf{B}}\) of B (III, § 3, no. 5, prop. 11), whence dim(B) = dim(A) (cor. 2 to prop. 8). This applies in particular to the following situation. * Let k be a complete non-discrete valued field, A the local ring of the polynomial ring \(k[\mathrm{X}_1, ..., \mathrm{X}_n]\) at the prime ideal generated by \(\mathrm{X}_1, ..., \mathrm{X}_n\) , and B the ring of convergent power series in \(\mathrm{X}_1, ..., \mathrm{X}_n\) with coefficients in k. Then the preceding hypotheses are satisfied, and consequently one has dim(B) = n. *
\end{enumerate}

\subsection{Construction of secant sequences}

PROPOSITION 9.--- Let A be a Noetherian ring, M a finitely generated A-module, and let a be a subset of A stable under addition and multiplication, \(\mathfrak{q}_{1},\ldots,\mathfrak{q}_{m}\) prime ideals of A not containing a, and \(r\geqslant1\) be an integer such that

\[
\operatorname{codim} (\mathrm{V} (\mathfrak {a}) \cap \operatorname{Supp} (\mathrm{M}), \operatorname{Supp} (\mathrm{M})) \geqslant r.\tag{21}
\]

There exists a sequence \((x_{1}, \ldots, x_{r})\) of elements of a, belonging to none of the ideals \(q_{1}, \ldots, q_{m}\) and such that the sequence \((x_{1}/1, \ldots, x_{r}/1)\) of elements of \(A_{p}\) is secant for the \(A_{p}\) -module \(M_{p}\) , for every prime ideal \(p \in V(a) \cap \text{Supp}(M)\) .

Let us reason by induction on \(r\). Suppose first that \(r = 1\), and denote by \(\Phi\) the (finite) set of minimal elements of \(\operatorname{Supp}(M)\). Argue by contradiction and suppose that there exists no element \(x_1\) of \(\mathfrak{a}\) satisfying the conditions of the statement. Let \(x \in \mathfrak{a}\), not belonging to \(\mathfrak{q}_1 \cup \ldots \cup \mathfrak{q}_m\). By hypothesis there exists a prime ideal \(\mathfrak{p}\) of \(V(\mathfrak{a}) \cap \operatorname{Supp}(M)\) such that the image \(x/1\) of \(x\) in \(A_{\mathfrak{p}}\) is not secant for \(M_{\mathfrak{p}}\). Let \(\Psi\) be the set of prime ideals \(\mathfrak{q} \in \Phi\) contained in \(\mathfrak{p}\); then the minimal elements of \(\operatorname{Supp}(M_{\mathfrak{p}})\) are the prime ideals \(\mathfrak{q}A_{\mathfrak{p}}\) of \(A_{\mathfrak{p}}\), where \(\mathfrak{q}\) ranges over \(\Psi\). By prop. 3 of no. 2, there therefore exists a \(\mathfrak{q} \in \Psi\) such that \(x/1 \in \mathfrak{q}A_{\mathfrak{p}}\), whence \(x \in \mathfrak{q}\). In other words, we have \(\mathfrak{a} \subset \mathfrak{q}_1 \cup \ldots \cup \mathfrak{q}_m \cup \bigcup_{\mathfrak{q} \in \Phi} \mathfrak{q}\). Since one has \(\mathfrak{a} \not\subset \mathfrak{q}_j\) for \(1 \leqslant j \leqslant m\), prop. 2 of II, § 1, no. 1 proves the existence of a prime ideal \(\mathfrak{q} \in \Phi\) containing \(\mathfrak{a}\), whence

\[
\mathrm{V} (\mathfrak {q}) \subset \mathrm{V} (\mathfrak {a}) \cap \operatorname{Supp} (\mathrm{M}).
\]

Since V(a) contains an irreducible component of Supp(M), this contradicts the hypothesis

\[
\operatorname{codim} (\mathrm{V} (\mathfrak {a}) \cap \operatorname{Supp} (\mathrm{M}), \operatorname{Supp} (\mathrm{M})) \geqslant 1.
\]

Suppose now \(r \geqslant 2\) . By the induction hypothesis, one can find a sequence \((x_1, ..., x_{r-1})\) of elements of \(\mathfrak{a}\) , not belonging to \(\mathfrak{q}_1 \cup ... \cup \mathfrak{q}_m\) , and such that, for every \(\mathfrak{p} \in V(\mathfrak{a}) \cap \operatorname{Supp}(M)\) , the sequence \((x_1/1, ..., x_{r-1}/1)\) of elements of \(A_{\mathfrak{p}}\) is secant for \(M_{\mathfrak{p}}\) . Put \(N = M / \sum_{i=1}^{r-1} x_i M\) . It suffices to construct an element \(x_r\) of \(\mathfrak{a}\) not belonging to \(\mathfrak{q}_1 \cup ... \cup \mathfrak{q}_m\) and such that, for every \(\mathfrak{p} \in V(\mathfrak{a}) \cap \operatorname{Supp}(M)\) , the element \(x_r/1\) of \(A_{\mathfrak{p}}\) is secant for \(N_{\mathfrak{p}}\) . By the first part of the proof, it suffices to establish the two relations

\begin{enumerate}
\def\labelenumi{(\arabic{enumi})}
\setcounter{enumi}{21}
\tightlist
\item
\end{enumerate}

\[
\mathrm{V} (\mathfrak {a}) \cap \operatorname{Supp} (\mathrm{M}) = \mathrm{V} (\mathfrak {a}) \cap \operatorname{Supp} (\mathrm{N}),\tag{23}
\]

\[
\operatorname{codim} (\mathrm{V} (\mathfrak {a}) \cap \operatorname{Supp} (\mathrm{N}), \operatorname{Supp} (\mathrm{N})) \geqslant 1.
\]

Now we have

\[
\operatorname{Supp} (\mathrm{N}) = \mathrm{V} (x _ {1}) \cap \dots \cap \mathrm{V} (x _ {r - 1}) \cap \operatorname{Supp} (\mathrm{M})
\]

by the corollary of prop. 18 of II, § 4, no. 4, and since \(x_1\) , \ldots, \(x_{r-1}\) belong to a, we deduce (22). Inequality (23) then follows from hypothesis (21) and from prop. 4, a) of no. 3, where one takes

\[
m = r - 1, \quad X = \operatorname{Supp} (M), \quad Y = V (\mathfrak {a}) \cap \operatorname{Supp} (N), \quad H _ {i} = V (x _ {i})
\]

whence X' = Supp(N).

COROLLARY 1. --- Let A be a Noetherian ring, M a finitely generated A-module, \(\mathfrak{p}_1, \ldots, \mathfrak{p}_n\) , \(\mathfrak{q}_1, \ldots, \mathfrak{q}_m\) prime ideals of A and \(r\) be an integer \(\geqslant 1\) . Suppose that one has \(\mathfrak{p}_i \notin \mathfrak{q}_j\) for \(1 \leqslant i \leqslant n\) and \(1 \leqslant j \leqslant m\) , and \(\dim_{\mathbf{A}_{\mathfrak{p}_i}} (\mathbf{M}_{\mathfrak{p}_i}) \geqslant r\) for \(1 \leqslant i \leqslant n\) . Then there exists a sequence \((x_1, \ldots, x_r)\) of elements of A belonging to all the \(\mathfrak{p}_i\) and to none of the \(\mathfrak{q}_j\) , such that, for \(1 \leqslant i \leqslant n\) , the images of \(x_1, \ldots, x_r\) in \(\mathbf{A}_{\mathfrak{p}_i}\) form a secant sequence for the module \(M_{\mathfrak{p}_i}\) .

Set \(\mathfrak{a} = \bigcap_{i}\mathfrak{p}_{i}\) . We have \(\mathbf{M}_{\mathfrak{p}_i}\neq \{0\}\) , hence \(\mathfrak{p}_i\in \mathrm{Supp}(\mathbf{M})\) for \(1\leqslant i\leqslant n\) , whence \(\mathrm{V}(\mathfrak{a})\subset \mathrm{Supp}(\mathbf{M})\) . We have

\[
\begin{array}{r l} \operatorname{codim} (\mathrm{V} (\mathfrak {a}) \cap \operatorname{Supp} (\mathrm{M}), \operatorname{Supp} (\mathrm{M})) & = \operatorname{codim} (\mathrm{V} (\mathfrak {a}), \operatorname{Supp} (\mathrm{M})) = \\ & = \inf _ {i} \left(\operatorname{codim} (\mathrm{V} (\mathfrak {p} _ {i}), \operatorname{Supp} (\mathrm{M}))\right) = \inf _ {i} \dim (\mathrm{M} _ {\mathfrak {p} _ {i}}) \geqslant r \end{array}
\]

(§ 1, no. 4, prop. 9), and one applies prop. 9.

COROLLARY 2. --- Let A be a Noetherian local ring, M a nonzero finitely generated A-module, and let a be a subset of \(\mathfrak{m}_{\mathbf{A}}\) , stable under addition and multiplication and such that long(M/aM) \textless{} + ∞. There exists a subset of a which is maximal secant for M.

Indeed, we have \(\mathrm{V}(\mathfrak{a}) \cap \operatorname{Supp}(\mathbf{M}) = \operatorname{Supp}(\mathbf{M}/\mathfrak{a}\mathbf{M}) = \{\mathfrak{m}_{\mathbf{A}}\}\) (§ 1, no. 4, remark 1), hence \(\operatorname{codim}(\mathrm{V}(\mathfrak{a}) \cap \operatorname{Supp}(\mathbf{M}), \operatorname{Supp}(\mathbf{M})) = \dim(\operatorname{Supp}(\mathbf{M})) = \dim_{\mathbf{A}}(\mathbf{M})\) , and one applies prop. 9.

\section{HILBERT-SAMUEL SERIES}

\subsection{The ring Z((T))}

Let A be a ring. Endow the A-module \(\mathbf{A}^{\mathbf{Z}}\) with the product topology of the discrete topologies. The elements \((a_{n})\in\mathbf{A}^{\mathbf{Z}}\) such that there exists \(n_{0}\in\mathbf{Z}\) with \(a_{n}=0\) for \(n<n_{0}\) form a submodule B of \(\mathbf{A}^{\mathbf{Z}}\) . If \(a=(a_{n})\in\mathbf{B}\), \(b=(b_{n})\in\mathbf{B}\), and we set \(ab=c\) , with \(c_{n}=\sum_{i+j=n}a_{i}b_{j}\) , one defines on B a structure of A-algebra. Let T be the element \((\theta_{n})\) of B such that \(\theta_{n}=0\) for \(n\neq1\) and \(\theta_{1}=1\) . Then T is invertible in B; for every element \(a=(a_{n})\) of B, the family \((a_{n}\mathrm{T}^{n})_{n\in\mathbf{Z}}\) is summable in \(\mathbf{A}^{\mathbf{Z}}\) and one has

\[
a = \sum_ {n \in \mathbf {Z}} a _ {n} \mathrm{T} ^ {n}.
\]

In the rest of this chapter, we denote by A((T)) the A-algebra B; it contains as subalgebras the algebra A{[}{[}T{]}{]} of formal series and the algebra A{[}T, T \(^{-1}\) {]}; their intersection is the algebra A{[}T{]} of polynomials.

Remark. --- The ring A((T)) is naturally identified with the ring of fractions A{[}{[}T{]}{]} \(_{T}\) of the ring A{[}{[}T{]}{]} defined by the multiplicative subset consisting of the powers of T.

For \(n, p\) in \(\mathbf{Z}\) , we define the integer \(\begin{bmatrix} n \\ p \end{bmatrix}\) by

\[
\left\{ \begin{array}{l} {\left[ \begin{array}{l} n \\ p \end{array} \right] = 0 \quad \text {if} \quad p <   0 \quad \text {or} \quad p > n,} \\ {\left[ \begin{array}{l} n \\ p \end{array} \right] = \binom {n} {p} = \frac {n (n - 1) \dots (n - p + 1)}{p !} \quad \text {if} \quad 0 \leqslant p \leqslant n.} \end{array} \right.\tag{1}
\]

We have \(\begin{bmatrix} n \\ p \end{bmatrix} = \begin{bmatrix} n \\ n - p \end{bmatrix}\) for \(n, p \in \mathbf{Z}\) .

Lemma 1.--- The element 1 - T of \(\mathbf{Z}((\mathrm{T}))\) is invertible. For every integer \(r > 0\) we have

\[
(1 - \mathrm{T}) ^ {- r} = \sum_ {n \in \mathbf {Z}} \left[ \begin{array}{c} n + r - 1 \\ r - 1 \end{array} \right] \mathrm{T} ^ {n} = \sum_ {n \in \mathbf {N}} \binom {n + r - 1} {r - 1} \mathrm{T} ^ {n}.
\]

Indeed, 1 - T is invertible in the ring \(\mathbf{Z}[[\mathrm{T}]]\) , with inverse \(\sum_{m\geqslant 0}\mathbf{T}^m\) ; we thus have

\[
(1 - \mathrm{T}) ^ {- r} = (\sum_ {m \geqslant 0} \mathrm{T} ^ {m}) ^ {r} = \sum_ {m _ {1}, \dots , m _ {r} \geqslant 0} \mathrm{T} ^ {m _ {1} + m _ {2} + \dots + m _ {r}},
\]

and the announced formula follows from E, III, p. 44, prop. 15.

Let \(\mathrm{Q}(\mathrm{T}) \in \mathbf{Z}[\mathrm{T}, \mathrm{T}^{-1}]\) , \(r\) be an integer \(> 0\) , and \(\mathrm{F} = (1 - \mathrm{T})^{-r}\mathrm{Q} \in \mathbf{Z}((\mathrm{T}))\) . Put

\[
\mathrm{Q} (\mathrm{T}) = \sum_ {i \in \mathbf {Z}} a _ {i} \mathrm{T} ^ {i}, \quad \mathrm{F} = \sum_ {n \in \mathbf {Z}} \alpha_ {n} \mathrm{T} ^ {n}.
\]

Then, by Lemma 1, we have

\[
\alpha_ {n} = \sum_ {i \in \mathbf {Z}} a _ {i} \left[ \begin{array}{c} n - i + r - 1 \\ r - 1 \end{array} \right] = \sum_ {i \leqslant n} a _ {i} \binom {n - i + r - 1} {r - 1}.
\]

Let \(n_1\) be the supremum in \(\overline{\mathbf{R}}\) of the set of integers \(i \in \mathbf{Z}\) such that \(a_i \neq 0\) . For every integer \(n \geqslant n_1\) , one has \(\alpha_n = \tilde{\alpha}(n)\) , where \(\tilde{\alpha}\) is the polynomial of \(\mathbf{Q}[\mathbf{X}]\) defined by

\[
\tilde {\alpha} (X) = \frac {1}{(r - 1) !} \sum_ {i \in \mathbf {Z}} a _ {i} \prod_ {j = 1} ^ {r - 1} (X - i + j).
\]

If we set \(c = \mathbf{Q}(1) = \sum_{i \in \mathbf{Z}} a_i\) , one has \(\tilde{\alpha}(\mathbf{X}) = c\mathbf{X}^{r-1}/(r - 1)! + \theta(\mathbf{X})\) , where \(\theta\) is a polynomial of degree \(\leqslant r - 2\) . Consequently, we have

\[
\alpha_ {n} = c \frac {n ^ {r - 1}}{(r - 1) !} + \rho_ {n} n ^ {r - 2},\tag{2}
\]

where the rational number \(\rho_{n}\) tends to a limit as \(n\) increases indefinitely. We deduce the relation

\[
\mathrm{Q} (1) = (r - 1)! \lim _ {n \rightarrow \infty} n ^ {1 - r} \alpha_ {n}.\tag{3}
\]

If \(F = \sum_{n \in \mathbf{Z}} a_n T^n\) and \(G = \sum_{n \in \mathbf{Z}} b_n T^n\) are two elements of \(Z((T))\) , we denote by “\(F \leqslant G\)” the relation “\(a_n \leqslant b_n\) for every \(n \in \mathbf{Z}\).” It is an order relation compatible with the ring structure of \(Z((T))\) (A, VI, p. 18, def. 1). We have \((1 - T)^{-1} \geqslant 1\) . If \(Q \in Z[T, T^{-1}]\) is \(\geqslant 0\) , then the integer \(Q(1)\) is positive.

Lemma 2. --- a) Let F be a nonzero element of \(\mathbf{Z}((\mathbf{T}))\) such that there exists \(r \in \mathbf{Z}\) , with \((1 - \mathrm{T})^r\mathrm{F} \in \mathbf{Z}[\mathrm{T}, \mathrm{T}^{-1}]\) ; then F is written uniquely in the form \(\mathrm{F} = (1 - \mathrm{T})^{-d}\) . Q, where \(\mathrm{Q} \in \mathbf{Z}[\mathrm{T}, \mathrm{T}^{-1}]\) , \(\mathrm{Q}(1) \neq 0\) and \(d \in \mathbf{Z}\) . If \(\mathrm{F} \geqslant 0\) , then we have \(\mathrm{Q}(1) > 0\) and \(d \geqslant 0\) .

\begin{enumerate}
\def\labelenumi{\alph{enumi})}
\setcounter{enumi}{1}
\tightlist
\item
  Let Q, R be in Z{[}T, T \(^{-1}\) {]}, and d, d' in Z with Q(1) \textgreater{} 0. If
\end{enumerate}

\[
(1 - \mathrm{T}) ^ {- d}. \mathrm{Q} \leqslant (1 - \mathrm{T}) ^ {- d ^ {\prime}}. \mathrm{R},
\]

then either \(d < d'\) , or \(d = d'\) and \(\mathbf{Q}(1) \leqslant \mathbf{R}(1)\) .

\begin{enumerate}
\def\labelenumi{\alph{enumi})}
\item
  One can write \(F = (1 - T)^{-r} T^{n} P(T)\) with r, \(n \in Z\) and \(P(T) \in Z[T]\) . By Euclidean division, one can write \(P(T) = (1 - T)^{p} R(T)\) with \(R(T) \in Z[T]\) and \(R(1) \neq 0\). Therefore \(F = (1 - T)^{-(r-p)} Q(T)\) , where \(Q(T) = T^{n} R(T) \in Z[T, T^{-1}]\) and \(Q(1) \neq 0\) . This proves the existence of d and Q. Moreover, if \((1 - T)^{r} Q(T) = (1 - T)^{s} R(T)\) with r \textgreater{} s and Q, R in \(Z[T, T^{-1}]\) , one has \(R(T) = (1 - T)^{r-s} Q(T)\) , hence \(R(1) = 0\) ; this proves uniqueness. Suppose that F is \(\geqslant 0\) ; if one had d \textless{} 0, then one would have \(F(1) = 0\) which is impossible since F is nonzero and all its coefficients are positive; we therefore have \(d \geqslant 0\) . If d = 0, then Q = F \(\geqslant 0\) , hence Q(1) is positive. If \(d \geqslant 1\) , then Q(1) is positive by formula (3). This proves a).
\item
  Suppose \(d \geqslant d'\) . Then \((1 - T)^{-d}((1 - T)^{d - d'} R - Q) \geqslant 0\) ; since \(S(T) = (1 - T)^{d - d'} R - Q\) belongs to \(Z[T, T^{-1}]\) , this implies \(S(1) \geqslant 0\) from what precedes. If \(d > d'\) , one has \(S(1) = -Q(1) < 0\) , whence a contradiction; if \(d = d'\) , one has \(S(1) = R(1) - Q(1)\) whence \(Q(1) \leqslant R(1)\) .
\end{enumerate}

\subsection{Poincaré series of a graded module over a polynomial ring}

Let \(H_{0}\) be a ring, I a finite set, and H the polynomial ring \(\mathrm{H}_{0}\big[(\mathrm{X}_{i})_{i\in\mathbf{I}}\big]\) . For each \(i\in I\) , let \(d_{i}\) be an integer \(>0\). Let us endow H with the structure of a graded ring of type Z such that the elements of \(H_{0}\) are homogeneous of degree 0 and each \(X_{i}\) is homogeneous of degree \(d_{i}\) . When \(d_{i}=1\) for every \(i\) , we recover the usual grading of polynomial rings.

Let M be a finitely generated graded H-module all of whose homogeneous components are \(H_{0}\) -modules of finite length; one calls the Poincaré series of M the element \(P_{M}\) of \(\mathbf{Z}((T))\) such that \(\mathrm{P}_{\mathrm{M}} = \sum_{n \in \mathbf{Z}} \mathrm{long}_{\mathrm{H}_{0}}(\mathrm{M}_{n}) \cdot \mathrm{T}^{n}\) , and one sets \(Q_{M} = P_{M} \cdot \prod_{i \in I} (1 - T^{d_{i}})\) .

THEOREM 1. --- The element \(\mathbf{Q}_{\mathbf{M}}\) of \(\mathbf{Z}((\mathbf{T}))\) belongs to \(\mathbf{Z}[\mathbf{T}, \mathbf{T}^{-1}]\) .

Dividing \(H_{0}\) by the annihilator of the \(H_{0}\) -module M, we reduce to the case where M is a faithful \(H_{0}\) -module. If \(a, b \in Z\) are such that M is generated as an H-module by \(M' = \sum_{a \leqslant i \leqslant b} M_{i}\) , then \(M'\) is a nonzero finitely generated \(H_{0}\) -module faithful and of finite length; consequently,

the ring \(\mathbf{H}_0\) is Artinian (A, VIII, § 1, no. 3), hence Noetherian (loc. cit., § 9, no. 1). The polynomial ring H is therefore Noetherian (loc. cit., § 1, no. 4). If I is empty, we have \(\mathbf{H} = \mathbf{H}_0\) , and the family of integers \((\mathrm{long}_{\mathbf{H}_0}(\mathbf{M}_n))_{n \in \mathbf{Z}}\) has finite support, since M is a finitely generated \(\mathbf{H}_0\) -module; whence the theorem in this case. We then argue by induction on the cardinality of the set I, assumed nonempty; let \(j \in \mathbf{I}\) and \(\mathbf{J} = \mathbf{I} - \{j\}\) . Denote by H' the graded subring of H generated by \(\mathbf{H}_0\) and the \(X_i\) for \(i\) in J; consider the homothety \((\mathbf{X}_j)_\mathbf{M}\) of ratio \(X_j\) on M, its kernel R, and its cokernel S. For each \(n \in \mathbf{Z}\) , we have an exact sequence of \(\mathbf{H}_0\) -modules

\[
0 \rightarrow \mathrm{R} _ {n - d _ {j}} \rightarrow \mathrm{M} _ {n - d _ {j}} \rightarrow \mathrm{M} _ {n} \rightarrow \mathrm{S} _ {n} \rightarrow 0,
\]

hence \(\mathbf{R}_{n - d_j}\) and \(\mathbf{S}_n\) are of finite length, and we have

\[
\operatorname{long} _ {\mathrm{H} _ {0}} \left(\mathrm{M} _ {n}\right) - \operatorname{long} _ {\mathrm{H} _ {0}} \left(\mathrm{M} _ {n - d _ {j}}\right) = \operatorname{long} _ {\mathrm{H} _ {0}} \left(\mathrm{S} _ {n}\right) - \operatorname{long} _ {\mathrm{H} _ {0}} \left(\mathrm{R} _ {n - d _ {j}}\right).\tag{4}
\]

Since M is a finitely generated module over the Noetherian ring H, the H-modules R and S are finitely generated; as they are annihilated by \(X_{j}\) , they are finitely generated H'-modules. By the induction hypothesis, the elements \(P_{R}\cdot\prod_{i\in J}(1-T^{d_{i}})\) and \(P_{S}\cdot\prod_{i\in J}(1-T^{d_{i}})\) of \(\mathbf{Z}((T))\) therefore belong to \(\mathbf{Z}[T,T^{-1}]\) ; from (4), we deduce

\[
\mathrm{P} _ {\mathrm{M}} - \mathrm{T} ^ {d _ {j}}. \mathrm{P} _ {\mathrm{M}} = \mathrm{P} _ {\mathrm{S}} - \mathrm{T} ^ {d _ {j}}. \mathrm{P} _ {\mathrm{R}},
\]

that is, \((1 - \mathrm{T}^{d_j}).\mathrm{P}_{\mathrm{M}} = \mathrm{P}_{\mathrm{S}} - \mathrm{T}^{d_j}.\mathrm{P}_{\mathrm{R}}\) ; we thus have

\[
\mathrm{P} _ {\mathrm{M}} \cdot \prod_ {i \in \mathrm{I}} (1 - \mathrm{T} ^ {d _ {i}}) = \mathrm{P} _ {\mathrm{S}} \cdot \prod_ {i \in \mathrm{J}} (1 - \mathrm{T} ^ {d _ {i}}) - \mathrm{T} ^ {d _ {j}} \cdot \mathrm{P} _ {\mathrm{R}} \cdot \prod_ {i \in \mathrm{J}} (1 - \mathrm{T} ^ {d _ {i}}),\tag{5}
\]

whence the conclusion.

Example 1. --- Suppose \(H_{0}\) is Artinian and take M = H. Then, with the previous notation, we have R = 0 and S = H', so by (5), we have \(Q_{H} = Q_{H'}\) ; since we have \(Q_{H_{0}} = \text{long}(H_{0})\) from which we obtain by induction \(Q_{H} = \text{long}(H_{0})\) , that is,

\[
\mathrm{P} _ {\mathrm{H}} = \operatorname{long} (\mathrm{H} _ {0}). \prod_ {i \in \mathrm{I}} (1 - \mathrm{T} ^ {d _ {i}}) ^ {- 1}.\tag{6}
\]

Suppose henceforth that H is equipped with the usual grading, for which \(d_{i}=1\) for every \(i\in I\) , and set \(r=\operatorname{Card}(\mathbf{I})\) ; one has \(\mathrm{P}_{\mathrm{M}}=\mathrm{Q}_{\mathrm{M}}(\mathrm{T}).(1-\mathrm{T})^{-r}\) . Put \(c_{\mathrm{M}}=\mathrm{Q}_{\mathrm{M}}(1)\) . Then, by formula (2) of no. 1, we have:

COROLLARY. --- a) If r = 0, then we have long \(_{H_{0}}(M)\) = c \(_{M}\) .

\begin{enumerate}
\def\labelenumi{\alph{enumi})}
\setcounter{enumi}{1}
\item
  If \(r = 1\) , then we have \(\mathrm{long}_{\mathrm{H}_0}(\mathbf{M}_n) = c_{\mathbf{M}}\) for \(n\) sufficiently large.
\item
  If \(r > 1\) , then we have long \(_{\mathrm{H}_0}(\mathbf{M}_n) = c_{\mathbf{M}} \frac{n^{r-1}}{(r - 1)!} + \rho_n n^{r-2}\) , where \(\rho_n\) tends to a limit in \(\mathbf{R}\) when \(n\) increases indefinitely.
\end{enumerate}

Remarks. --- 1) The integer \(c_{\mathbf{M}}\) is positive by Lemma 2. One can have \(\mathbf{M} \neq 0\) and \(c_{\mathbf{M}} = 0\) (cf. prop. 2).

\begin{enumerate}
\def\labelenumi{\arabic{enumi})}
\setcounter{enumi}{1}
\tightlist
\item
  Let \(0 \to M' \to M \to M'' \to 0\) be an exact sequence of graded H-modules and degree-0 homomorphisms such that M is finitely generated over H and \(M_n\) is of finite length over \(H_0\) for each \(n\). Then, for each \(n \in \mathbf{Z}\) , one has
\end{enumerate}

\[
\operatorname{long} _ {\mathrm{H} _ {0}} \left(\mathrm{M} _ {n}\right) = \operatorname{long} _ {\mathrm{H} _ {0}} \left(\mathrm{M} _ {n} ^ {\prime}\right) + \operatorname{long} _ {\mathrm{H} _ {0}} \left(\mathrm{M} _ {n} ^ {\prime \prime}\right),
\]

hence \(\mathrm{P_M} = \mathrm{P_{M'}} + \mathrm{P_{M''}}\) , \(\mathrm{Q_M} = \mathrm{Q_{M'}} + \mathrm{Q_{M''}}\) and \(c_{\mathbf{M}} = c_{\mathbf{M'}} + c_{\mathbf{M''}}\) .

\begin{enumerate}
\def\labelenumi{\arabic{enumi})}
\setcounter{enumi}{2}
\tightlist
\item
  Let \(\mathbf{M}(p)\) be the module deduced from \(\mathbf{M}\) by shift by \(p\) in the grading (A, II, p. 165, example 3). Since we have \(\mathbf{M}(p)_n = \mathbf{M}_{p+n}\) , one has \(\mathbf{P}_{\mathbf{M}(p)} = \mathrm{T}^{-p}\mathbf{P}_{\mathbf{M}}\) , \(\mathbf{Q}_{\mathbf{M}(p)} = \mathrm{T}^{-p}\mathbf{Q}_{\mathbf{M}}\) and \(c_{\mathbf{M}(p)} = c_{\mathbf{M}}\) .
\end{enumerate}

Examples. --- 2) Suppose H \(_{0}\) is Artinian, and let M be a graded free H-module generated by s homogeneous, linearly independent elements of respective degrees \(\delta_{1}, ..., \delta_{s}\). Then M is isomorphic to H( \(-\delta_{1}\) ) ⊕ \(\cdots\) ⊕ H( \(-\delta_{s}\) ). By Remarks 2 and 3 and Example 1, we therefore have

\[
\begin{array}{l} \mathrm {P_ {M}} = \operatorname{long} (\mathrm {H_ {0}}) \left(\sum_ {i = 1} ^ {s} \mathrm{T} ^ {\delta_ {i}}\right) (1 - \mathrm{T}) ^ {- r}, \\ \mathrm {Q_ {M}} = \operatorname{long} (\mathrm {H_ {0}}) \left(\sum_ {i = 1} ^ {s} \mathrm{T} ^ {\delta_ {i}}\right), \\ c _ {\mathrm{M}} = s. \operatorname{long} (\mathrm {H_ {0}}). \end{array}
\]

\begin{enumerate}
\def\labelenumi{\arabic{enumi})}
\setcounter{enumi}{2}
\tightlist
\item
  Suppose again that \(H_{0}\) is Artinian; let M be a graded H-module, and suppose there exists an exact sequence of graded H-modules and degree-0 homomorphisms
\end{enumerate}

\[
0 \to \mathrm{L} _ {n} \to \mathrm{L} _ {n - 1} \to \dots \to \mathrm{L} _ {0} \to \mathrm{M} \to 0,
\]

where, for \(k=0,1,...,n\) , \(\mathbf{L}_{k}\) is a free graded H-module generated by linearly independent homogeneous elements, of respective degrees \(\delta_{k,1},..., \delta_{k,m(k)}\). Then, by Remark 2 and Example 2, we have

\[
\mathrm{Q} _ {\mathrm{M}} = \operatorname{long} (\mathrm{H} _ {0}). \sum_ {0 \leqslant k \leqslant n} \sum_ {1 \leqslant j \leqslant m (k)} (- 1) ^ {k} \mathrm{T} ^ {\delta_ {k, j}},
\]

\[
c _ {\mathbf {M}} = \operatorname{long} (\mathrm{H} _ {0}). \sum_ {0 \leqslant k \leqslant n} (- 1) ^ {k} m (k).
\]

Remark 4. --- One can prove (p. 88, exerc. 4) that under the hypotheses of th. 1, the \(\mathrm{H}_{0}\) -modules \(\mathrm{Tor}_{j}^{\mathrm{H}}(\mathrm{H}_{0},\mathbf{M})\) are of finite length, are zero for \(j > r\) , and one has \(c_{\mathbf{M}} = \sum_{j=0}^{r} (-1)^{j} \mathrm{long}_{\mathrm{H}_{0}}(\mathrm{Tor}_{j}^{\mathrm{H}}(\mathrm{H}_{0},\mathbf{M}))\). More precisely, the H-modules

\(T_{j} = \text{Tor}_{j}^{\text{H}}(\text{H}_{0}, \text{M})\) are naturally equipped with graduations, and one has

\[
\mathrm{Q} _ {\mathrm{M}} = \sum_ {j = 0} ^ {r} (- 1) ^ {j} \mathrm{P} _ {\mathrm{T} _ {j}}.
\]

PROPOSITION 1.--- Let M be a graded H-module. Suppose that M is generated by \(M_{0}\) and that \(M_{0}\) is an \(H_{0}\)-module of finite length. Then we have

\[
\mathrm{P} _ {\mathrm{M}} \leqslant (1 - \mathrm{T}) ^ {- r} \operatorname{long} _ {\mathrm{H} _ {0}} (\mathrm{M} _ {0}), \quad c _ {\mathrm{M}} \leqslant \operatorname{long} _ {\mathrm{H} _ {0}} (\mathrm{M} _ {0}).
\]

Moreover, the following conditions are equivalent:

\begin{enumerate}
\def\labelenumi{(\roman{enumi})}
\tightlist
\item
  \(c_{\mathbf{M}} = \mathrm{long}_{\mathrm{H_0}}(\mathbf{M}_0)\) ;
\end{enumerate}

\[
\text{(ii)} \mathrm{P}_{\mathbf{M}} = \operatorname{long}_{\mathbf{H}_{0}}(\mathbf{M}_{0}) \cdot (1 - \mathrm{T})^{-r},\ \text{that is, }\mathbf{M} = \mathbf{M}_{0}\text{ if }r=0\text{ and}
\]

\[
\operatorname{long} _ {\mathrm{H} _ {0}} \left(\mathrm{M} _ {n}\right) = \operatorname{long} _ {\mathrm{H} _ {0}} \left(\mathrm{M} _ {0}\right). \binom {n + r - 1} {r - 1}
\]

for \(n \in \mathbf{N}\) if \(r > 0\) ;

\begin{enumerate}
\def\labelenumi{(\roman{enumi})}
\setcounter{enumi}{2}
\tightlist
\item
  the canonical homomorphism of H-modules
\end{enumerate}

\[
\varphi : \mathrm{H} \otimes_ {\mathrm{H} _ {0}} \mathrm{M} _ {0} \rightarrow \mathrm{M}
\]

is bijective.

Let R denote the kernel of φ. Since φ is surjective, we have

\(P_{M} = P_{H \otimes M_{0}} - P_{R} = \text{long}_{H_{0}}(M_{0})(1 - T)^{-r} - P_{R}\) and \(c_{M} = \text{long}_{H_{0}}(M_{0}) - c_{R}\). The conditions (i), (ii), and (iii) are respectively equivalent to \(c_{R} = 0\), \(P_{R} = 0\) and \(R = 0\). We therefore have (iii) \(\Rightarrow\) (ii) \(\Rightarrow\) (i), and it suffices to prove that \(c_{R} = 0\) implies \(R = 0\). Suppose \(R \neq 0\) and let \(0 = N^{h} \subset N^{h-1} \subset \ldots \subset N^{0} = M_{0}\) be a Jordan–Hölder sequence of the \(H_{0}\)-module \(M_{0}\). Let \(R^{m}\) be the intersection of \(R\) and the image of \(H \otimes_{H_{0}} N^{m}\) in \(H \otimes_{H_{0}} M_{0}\); there exists an integer \(m\) between 1 and \(h\) such that \(R^{m} \neq R^{m-1}\). Put \(L = R^{m-1}/R^{m}\); one has \(0 \leqslant c_{L} \leqslant c_{R}\) and it suffices to prove that \(c_{L} \neq 0\). Now, if \(k\) is the quotient field of \(H_{0}\) by the maximal annihilator ideal of \(N^{m-1}/N^{m}\), L is identified with a nonzero graded submodule of \(k[(X_{i})_{i\in I}]\). Therefore L contains a submodule isomorphic to a shifted module of \(k[(X_{i})_{i\in I}]\); since \(c_{k[(X_{i})_{i\in I}]} = 1\), one therefore has \(c_{L} \geqslant 1\) (remarks 2 and 3), which is what we wanted to prove.

By A, X, p. 160, th. 1, condition (iii) means that \((\mathbf{X}_{1}, \ldots, \mathbf{X}_{r})\) is a completely secant sequence for the H-module M.

PROPOSITION 2. --- Suppose that \(H_{0}\) is a field, and let M be a finitely generated graded H-module. Let K be the field of fractions of H. Then \(c_{M}\) is equal to the rank of the H-module M, that is, to the dimension of the K-vector space \(M \otimes_{H} K\) .

This is clear if M = H, since \(c_{H} = 1\) . Moreover, let \(x \in H\) , homogeneous of degree d, and nonzero; we have (H/xH) \(\otimes_{H} K = 0\) ; from the exact sequence

\[
0 \rightarrow \mathrm{H} (- d) \rightarrow \mathrm{H} \rightarrow \mathrm{H} / x \mathrm{H} \rightarrow 0,
\]

and Remarks 2 and 3, we deduce \(c_{\mathrm{H}/x\mathrm{H}}=0\) . The proposition is therefore verified when M is generated by a single homogeneous element. The general case follows from this, since every finitely generated graded H-module has a composition series whose quotients are of the preceding form.

Remark 6. -- Under the hypotheses of Prop. 2, we therefore have \(c_{\mathbf{M}} = 0\) if and only if M is a torsion H-module, or equivalently if and only if \(\dim_{\mathrm{H}}(\mathbf{M}) < r\) ( \(\S 1, \mathfrak{n}^{\circ} 5\) , Example 4).

\subsection{Hilbert-Samuel series of a well-filtered module}

In the remainder of this paragraph, we will use the following notation: if \(\mathbf{G} \in \mathbf{Z}((\mathbf{T}))\) and if \(r \in \mathbf{N}\) , we set \(\mathbf{G}^{(r)} = (1 - \mathbf{T})^{-r}\mathbf{G}\) ; in particular, if \(\mathbf{G} = \sum_{n \in \mathbf{Z}} a_n \mathbf{T}^n\) , then

\[
\mathrm{G} ^ {(1)} = \sum_ {n \in \mathbf {Z}} \left(\sum_ {i \leqslant n} a _ {i}\right) \mathrm{T} ^ {n}.
\]

If \(\mathbf{G} \geqslant 0\) , one has \(\mathbf{G}^{(r)} \geqslant 0\) for every \(r \in \mathbf{N}\) .

Let A be a Noetherian ring, q an ideal of A, and M a finitely generated A-module. Recall (III, § 3, no. 1, def. 1) that a q-good filtration on M is a map F : n ↦ Fn from Z to the set of submodules of M satisfying the following three conditions:

\begin{enumerate}
\def\labelenumi{\alph{enumi})}
\item
  we have \(\mathfrak{q}F_n\subset \mathrm{F}_{n + 1}\subset \mathrm{F}_n\) for every \(n\in \mathbf{Z},\)
\item
  there exists \(n_0 \in \mathbf{Z}\) such that \(\mathfrak{q}F_n = \mathrm{F}_{n+1}\) for \(n \geqslant n_0\) ,
\item
  there exists \(n_{1} \in \mathbf{Z}\) such that \(\mathrm{F}_{n_1} = \mathbf{M}\) .
\end{enumerate}

If \(n_{0}\) and \(n_{1}\) satisfy the preceding conditions, we have, for all \(n\in\mathbf{Z}\) ,

\[
\mathfrak {q} ^ {n - n _ {1}} \mathrm{M} \subset \mathrm{F} _ {n} \subset \mathfrak {q} ^ {n - n _ {0}} \mathrm{M}\tag{7}
\]

(recall that we have set \(\mathfrak{q}^r = \mathbf{A}\) for \(r\leqslant 0\) , by convention).

Lemma 3.--- If F and F' are two q-good filtrations on M, there exists an integer m such that \(F_{n}^{\prime} \subset F_{n-m}\) for every \(n \in Z\) .

Indeed, there exists \(n_{2}\) such that \(\mathbf{F}_{n}^{\prime}\subset\mathfrak{q}^{n-n_{2}}\mathbf{M}\) for every \(n\) , hence \(\mathbf{F}_{n}^{\prime}\subset\mathbf{F}_{n-(n_{2}-n_{1})}\) for every \(n\) .

Lemma 4. --- Let F be a q-good filtration on M. If M/qM has finite length, then M/F \(_{n+1}\) and F \(_{n}\) /F \(_{n+1}\) have finite length for all n ∈ Z.

With the notation from (7), we have long(M/F \(_{n+1}\) ) ≤ long(M/q \(^{n-n_1+1}\) M), and it suffices to prove that q \(^{n}\) M/q \(^{n+1}\) M has finite length for every n. We are therefore reduced to the case of the q-adic filtration. Let (x \(_{1}\) , \ldots, x \(_{r}\) ) be a finite generating system of the A-module q, and let I be the finite set of monomials of total degree n in r variables

\(X_{1},...,X_{r}\) . The homomorphism from \((\mathbf{M}/\mathfrak{q}\mathbf{M})^{\mathbf{I}}\) into \(q^{n}M/q^{n+1}M\) which maps a family \((u_{m})_{m\in\mathbf{I}}\) to the element \(\sum m(x_{1},...,x_{r})u_{m}\) is surjective. Since M/qM has finite length, \(q^{n}M/q^{n+1}M\) also has finite length.

Assume now that M/qM has finite length. Let F be a q-good filtration on M. There exists \(n_{1} \in Z\) such that \(F_{n_{1}} = M\) , hence \(F_{n} = M\) for \(n \leqslant n_{1}\) ; we therefore define an element \(H_{M,F}\) of \(\mathbf{Z}((T))\) by setting

\[
\mathrm{H} _ {\mathrm{M}, \mathrm{F}} = \sum_ {n \in \mathbf {Z}} \operatorname{long} _ {\mathrm{A} / \mathfrak {q}} (\mathrm{F} _ {n} / \mathrm{F} _ {n + 1}). \mathrm{T} ^ {n} \in \mathbf {Z} ((\mathrm{T})) .\tag{8}
\]

DEFINITION 1. --- We call H \(_{M,F}\) the Hilbert-Samuel series of the A-module M (relative to the q-good filtration F).

The map \(n\mapsto\text{long}_{\mathbf{A}}(\mathbf{F}_{n}/\mathbf{F}_{n+1})\) is often called the Hilbert-Samuel function of M (relative to F).

This applies in particular to the case of the q-adic filtration (\(F_{n} = q^{n}M\)); we then set \(H_{M,F} = H_{M,q}\) . One therefore has

\[
\mathrm{H} _ {\mathrm{M}, \mathfrak {q}} = \sum_ {n \in \mathbf {Z}} \operatorname{long} _ {\mathrm{A} / \mathfrak {q}} (\mathfrak {q} ^ {n} \mathrm{M} / \mathfrak {q} ^ {n + 1} \mathrm{M}). \mathrm{T} ^ {n}.\tag{9}
\]

PROPOSITION 3.--- a) If F is a q-good filtration on M, then we have

\[
\mathrm{H} _ {\mathrm{M}, \mathrm{F}} ^ {(1)} = \sum_ {n \in \mathbf {Z}} \operatorname{long} _ {\mathrm{A}} (\mathrm{M} / \mathrm{F} _ {n + 1}). \mathrm{T} ^ {n}.\tag{10}
\]

\begin{enumerate}
\def\labelenumi{\alph{enumi})}
\setcounter{enumi}{1}
\tightlist
\item
  If F and F' are two q-good filtrations on M, there exists an integer m such that \(\mathrm{H}_{\mathbf{M},\mathrm{F}'}^{(1)} \geqslant \mathrm{T}^{m}\mathrm{H}_{\mathbf{M},\mathbf{F}}^{(1)}\) .
\end{enumerate}

The part \(a)\) follows immediately from the definition of \(\mathbf{H}_{\mathbf{M},\mathbf{F}}^{(1)}\) ; the part \(b)\) follows from \(a)\) and Lemma 3.

THEOREM 2. --- Let A be a Noetherian ring, q an ideal of A, M a finitely generated A-module such that M/qM is nonzero and of finite length, and F a q-good filtration on M.

\begin{enumerate}
\def\labelenumi{\alph{enumi})}
\item
  There exists an integer \(d \geqslant 0\) and an element \(\mathbf{R}\) of \(\mathbf{Z}[\mathrm{T}, \mathrm{T}^{-1}]\) , uniquely determined, such that \(\mathbf{R}(1) > 0\) and \(\mathrm{H}_{\mathrm{M,F}} = (1 - \mathrm{T})^{-d}\mathbf{R}\) .
\item
  The integers d and R(1) are independent of the chosen q-good filtration F.
\item
  Consider the graded ring \(\operatorname{gr}(A)\) such that \(\operatorname{gr}_{n}(A) = \mathfrak{q}^{n}/\mathfrak{q}^{n+1}\) and the graded \(\operatorname{gr}(A)\)-module \(\operatorname{gr}(M)\) such that \(\operatorname{gr}_{n}(M) = F_{n}/F_{n+1}\). Since we have \(F_{n_{1}} = M\) and \(\mathfrak{q}F_{n} = F_{n+1}\) for \(n \geqslant n_{0}\), \(\operatorname{gr}(M)\) is generated by \(\bigoplus_{n_{1} \leqslant n \leqslant n_{0}} \operatorname{gr}_{n}(M)\), hence is of finite type. Moreover, if \((x_{1},\ldots,x_{r})\) is a finite generating system of the A-module \(\mathfrak{q}\), \(\operatorname{gr}(A)\) is generated by \(\operatorname{gr}_{0}(A)\) and the classes of the \(x_{i}\) modulo \(\mathfrak{q}^{2}\), hence is isomorphic to a graded quotient ring of \(H = (A/\mathfrak{q})[X_{1},\ldots,X_{r}]\). By th. 1 of no. 2, we have
\end{enumerate}

\[
(1 - \mathrm{T}) ^ {r} \mathrm{H} _ {\mathrm{M}, \mathrm{F}} \in \mathbf {Z} [ \mathrm{T}, \mathrm{T} ^ {- 1} ].
\]

We have \(H_{M,F} \neq 0\) and there therefore exist \(d \in N\) and \(R \in Z[T, T^{-1}]\) uniquely determined such that \(\mathrm{R}(1) > 0\) and \(\mathrm{H}_{\mathrm{M},\mathrm{F}} = (1 - \mathrm{T})^{-d}. R\) (Lemma 2 of no. 1).

\begin{enumerate}
\def\labelenumi{\alph{enumi})}
\setcounter{enumi}{1}
\tightlist
\item
  Let F' be another q-good filtration and write similarly
\end{enumerate}

\[
\mathrm{H} _ {\mathbf {M}, \mathrm{F} ^ {\prime}} = (1 - \mathrm{T}) ^ {- d ^ {\prime}} \mathbf {R} ^ {\prime}.
\]

By Prop. 3, b), there exists an integer m such that \((1 - T)^{-d'-1}R' \geqslant T^{m}(1 - T)^{-d-1}R\) . By Lemma 2, b) of no. 1, this implies \(d' \geqslant d\) and, if \(d' = d\) , \(R'(1) \geqslant R(1)\) . Exchanging the roles of F and \(F'\) , we obtain \(d = d'\) and \(R(1) = R'(1)\) .

Remark 1. --- With the notations of \(a\) ), write \(\mathbf{R} = \sum_{i \in \mathbf{Z}} a_i \mathbf{T}^i\) , and suppose \(d > 0\) . By no. 1, the relation \(\mathrm{H}_{\mathrm{M},\mathrm{F}} = (1 - \mathrm{T})^{-d}\mathrm{R}\) is also written

\[
\operatorname{long} _ {\mathbf {A}} \left(\mathrm{F} _ {n} / \mathrm{F} _ {n + 1}\right) = \sum_ {i \in \mathbf {Z}} a _ {i} \left[ \begin{array}{c} n - i + d - 1 \\ d - 1 \end{array} \right] = \sum_ {i \leqslant n} a _ {i} \binom {n - i + d - 1} {d - 1}.\tag{11}
\]

Similarly, since \(\mathbf{H}_{\mathbf{M},\mathbf{F}}^{(1)} = (1 - \mathbf{T})^{-d - 1}\mathbf{R}\) , one has

\[
\operatorname{long} _ {\mathbf {A}} \left(\mathrm{M} / \mathrm{F} _ {n + 1}\right) = \sum_ {i \in \mathbf {Z}} a _ {i} \left[ \begin{array}{c} n - i + d \\ d \end{array} \right] = \sum_ {i \leqslant n} a _ {i} \binom {n - i + d} {d}.\tag{12}
\]

Let A be a Noetherian ring, q an ideal of A, and M a finitely generated A-module such that M/qM has finite length. If M ≠ qM, then by Theorem 2, b), there exist integers \(d_{q}(M) \geqslant 0\) and \(e_{q}(M) > 0\) such that, for every q-good filtration F on M, there exists R ∈ Z{[}T, T \(^{-1}\) {]} with

\[
\mathrm{H} _ {\mathbf {M}, \mathrm{F}} = (1 - \mathrm{T}) ^ {- d _ {\mathfrak {q}} (\mathbf {M})} \mathrm{R}, \quad \mathrm{R} (1) = e _ {\mathfrak {q}} (\mathbf {M}).
\]

If M = qM, we set by convention \(d_{q}(M) = -\infty\) , \(e_{q}(M) = 0\) .

Remark 2. --- Saying that M/qM has finite length means that

\[
\operatorname{Supp} (\mathrm{M} / \mathfrak {q} \mathrm{M}) = \operatorname{Supp} (\mathrm{M}) \cap \mathrm{V} (\mathfrak {q})
\]

is formed of maximal ideals (IV, § 2, no. 5, prop. 7). We shall see below (no. 4, corollary to th. 3) that \(d_{\mathfrak{q}}(\mathbf{M})\) is the least upper bound of the numbers \(\dim_{\mathbf{A}_{\mathfrak{m}}}(M_{\mathfrak{m}})\) where m ranges over the set \(\operatorname{Supp}(\mathbf{M}) \cap \mathbf{V}(\mathfrak{q})\) .

COROLLARY. — Let A be a Noetherian ring, q an ideal of A, M a finitely generated A-module such that M/qM has finite length, and F a q-good filtration on M.

\begin{enumerate}
\def\labelenumi{\alph{enumi})}
\tightlist
\item
  For one to have \(d_{\mathfrak{q}}(\mathbf{M}) \leqslant 0\) it is necessary and sufficient that the sequence \((\mathfrak{q}^{n}\mathbf{M})\) be stationary, or equivalently that the sequence \((\mathrm{F}_{n})\) be stationary. Then, for all sufficiently large n,
\end{enumerate}

\[
\operatorname{long} \left(\mathrm{M} / \mathrm{F} _ {n + 1}\right) = \operatorname{long} \left(\mathrm{M} / \mathfrak {q} ^ {n + 1} \mathrm{M}\right) = e _ {\mathfrak {q}} (\mathrm{M}).
\]

\begin{enumerate}
\def\labelenumi{\alph{enumi})}
\setcounter{enumi}{1}
\tightlist
\item
  Suppose that we have \(d_{\mathfrak{q}}(\mathbf{M}) > 0\) . Then one has
\end{enumerate}

\begin{enumerate}
\def\labelenumi{(\arabic{enumi})}
\setcounter{enumi}{12}
\tightlist
\item
\end{enumerate}

\[
\operatorname{long} _ {\mathbf {A}} \left(\mathrm{F} _ {n} / \mathrm{F} _ {n + 1}\right) = e _ {\mathfrak {q}} (\mathrm{M}) n ^ {d _ {\mathfrak {q}} (\mathrm{M}) - 1} / \left(d _ {\mathfrak {q}} (\mathrm{M}) - 1\right)! + \rho_ {n} n ^ {d _ {\mathfrak {q}} (\mathrm{M}) - 2},\tag{14}
\]

\[
\operatorname{long} _ {\mathbf {A}} (\mathrm{M} / \mathrm{F} _ {n + 1}) = e _ {\mathfrak {q}} (\mathbf {M}) n ^ {d _ {\mathfrak {q}} (\mathbf {M})} / d _ {\mathfrak {q}} (\mathbf {M})! + \sigma_ {n} n ^ {d _ {\mathfrak {q}} (\mathbf {M}) - 1},
\]

where \(\rho_{n}\) and \(\sigma_{n}\) tend to a limit as \(n\) increases indefinitely.

This follows immediately from th. 2 and formula (2) of no. 1.

Remark 3. --- Suppose q is contained in the radical of A. Then, by Nakayama's lemma, the sequence (q \(^{n}\) M) is stationary if and only if one has q \(^{n}\) M = 0 for n sufficiently large. It then follows from part a) of the corollary that one has d \(_{q}\) (M) ≤ 0 if and only if M is of finite length, and one then has e \(_{q}\) (M) = long \(_{A}\) (M).

PROPOSITION 4. — Let A be a Noetherian ring, \(x_{1},...,x_{r}\) elements of A, \(\mathfrak{x}\) the ideal they generate, and M a finitely generated A-module such that \(M/\mathfrak{x}M\) is nonzero and of finite length.

\begin{enumerate}
\def\labelenumi{\alph{enumi})}
\item
  We have \(d_{\mathfrak{x}}(\mathbf{M}) \leqslant r\) .
\item
  If \(d_{\mathfrak{x}}(M) = r\), then \(e_{\mathfrak{x}}(M) \leqslant \text{long}_{A}(M/\mathfrak{x}M)\).
\item
  If the sequence \((x_{1},...,x_{r})\) is completely secant for M (A, X, p. 157), then \(d_{\mathfrak{x}}(\mathbf{M}) = r\) and \(e_{\mathfrak{x}}(\mathbf{M}) = \text{long}_{\mathbf{A}}(\mathbf{M}/\mathfrak{x}\mathbf{M})\). The converse is true if the \(x_{i}\) belong to the radical of A.
\end{enumerate}

Let H be the polynomial ring \((\mathbf{A} / \mathfrak{x})[\mathbf{X}_1,\dots ,\mathbf{X}_r]\). We equip \(\mathbf{G} = \bigoplus_{n}\mathfrak{x}^{n}\mathbf{M} / \mathfrak{x}^{n + 1}\mathbf{M}\)

of the structure of graded H-module for which \((\mathrm{X}_{i})_{\mathrm{G}}\) is multiplication by the class of \(x_{i}\) modulo \(\mathfrak{x}^{2}\). With the notations \(P_{G}, Q_{G}, c_{G}\) from No. 2, we have \(H_{M,\mathfrak{x}} = P_{G}\) , hence \((1 - T)^{-d_{\mathfrak{x}}(M)}R = (1 - T)^{-r}Q_{G}\) , where \(R(1) = e_{\mathfrak{x}}(M) > 0\) and \(Q_{G}(1) = c_{G}\) .

Thus either \(d_{\mathfrak{x}}(\mathbf{M}) < r\) and \(c_{\mathbf{G}} = 0\), or \(d_{\mathfrak{x}}(\mathbf{M}) = r\) and \(c_{\mathbf{G}} = e_{\mathfrak{x}}(\mathbf{M})\). Moreover, by prop. 1 of no. 2, we have \(c_{\mathbf{G}} \leqslant \text{long}(\mathbf{M}/\mathfrak{x}\mathbf{M})\), and there is equality if and only if the canonical homomorphism \(\mathrm{A}/\mathfrak{x}[\mathrm{X}_{1}, ..., \mathrm{X}_{r}] \otimes_{\mathrm{A}/\mathfrak{x}} \mathrm{M}/\mathfrak{x}\mathrm{M} \to \bigoplus_{n}\mathfrak{x}^{n}\mathrm{M}/\mathfrak{x}^{n+1}\mathrm{M}\) is bijective.

This entails the proposition, in view of A, X, p. 160, th. 1.

PROPOSITION 5. --- Let \(0 \to \mathbf{M}' \to \mathbf{M} \to \mathbf{M}'' \to 0\) be an exact sequence of finitely generated modules over a Noetherian ring A, and q an ideal of A.

\begin{enumerate}
\def\labelenumi{\alph{enumi})}
\item
  For \(M/\mathfrak{q}M\) to have finite length, it is necessary and sufficient that the same be true of \(M'/\mathfrak{q}M'\) and \(M''/\mathfrak{q}M''\).
\item
  Suppose M/qM has finite length. Then we are in one of the following three cases:
\end{enumerate}

\[
1) d _ {\mathfrak {q}} (\mathbf {M}) = d _ {\mathfrak {q}} \left(\mathbf {M} ^ {\prime}\right) > d _ {\mathfrak {q}} \left(\mathbf {M} ^ {\prime \prime}\right) and e _ {\mathfrak {q}} (\mathbf {M}) = e _ {\mathfrak {q}} \left(\mathbf {M} ^ {\prime}\right),
\]

\begin{enumerate}
\def\labelenumi{\arabic{enumi})}
\setcounter{enumi}{1}
\tightlist
\item
  \(d_{\mathfrak{q}}(\mathbf{M}) = d_{\mathfrak{q}}(\mathbf{M}^{\prime \prime}) > d_{\mathfrak{q}}(\mathbf{M}^{\prime})\) and \(e_{\mathfrak{q}}(\mathbf{M}) = e_{\mathfrak{q}}(\mathbf{M}^{\prime \prime}),\)
\end{enumerate}

\[
3) d _ {\mathfrak {q}} (\mathbf {M}) = d _ {\mathfrak {q}} \left(\mathbf {M} ^ {\prime}\right) = d _ {\mathfrak {q}} \left(\mathbf {M} ^ {\prime \prime}\right) and e _ {\mathfrak {q}} (\mathbf {M}) = e _ {\mathfrak {q}} \left(\mathbf {M} ^ {\prime}\right) + e _ {\mathfrak {q}} \left(\mathbf {M} ^ {\prime \prime}\right).
\]

\begin{enumerate}
\def\labelenumi{\alph{enumi})}
\item
  We have \(\operatorname{Supp}(\mathbf{M}) = \operatorname{Supp}(\mathbf{M}') \cup \operatorname{Supp}(\mathbf{M}'')\) and the assertion follows from Remark 2.
\item
  Equip M with a q-good filtration F (for example, the q-adic filtration), M'' with the quotient filtration F'', and M' with the induced filtration F'. The filtrations F' and F'' are q-good (III, § 3, no. 1, prop. 1). Then for each n we have an exact sequence of A-modules
\end{enumerate}

\[
0 \rightarrow \mathrm{F} _ {n} ^ {\prime} / \mathrm{F} _ {n + 1} ^ {\prime} \rightarrow \mathrm{F} _ {n} / \mathrm{F} _ {n + 1} \rightarrow \mathrm{F} _ {n} ^ {\prime \prime} / \mathrm{F} _ {n + 1} ^ {\prime \prime} \rightarrow 0
\]

(III, § 2, no. 4, prop. 2), so that we have \(\mathrm{H}_{\mathbf{M},\mathbf{F}} = \mathrm{H}_{\mathbf{M}',\mathbf{F}'} + \mathrm{H}_{\mathbf{M}'',\mathbf{F}''}\) , or equivalently

\[
(1 - \mathrm{T}) ^ {- d _ {\mathfrak {q}} (\mathrm{M})} \mathrm{R} = (1 - \mathrm{T}) ^ {- d _ {\mathfrak {q}} (\mathrm{M} ^ {\prime})} \mathrm{R} ^ {\prime} + (1 - \mathrm{T}) ^ {- d _ {\mathfrak {q}} (\mathrm{M} ^ {\prime \prime})} \mathrm{R} ^ {\prime \prime},
\]

with \(R,R',R''\in\mathbf{Z}[T,T^{-1}]\), \(R(1)=e_{\mathfrak{q}}(\mathbf{M})\), \(R'(1)=e_{\mathfrak{q}}(\mathbf{M}')\), \(R''(1)=e_{\mathfrak{q}}(\mathbf{M}'')\). Assertion b) follows immediately from this.

\subsection{Degree of the Hilbert-Samuel function}

THEOREM 3. --- Let A be a Noetherian local ring, q an ideal of A distinct from A, and M a finitely generated A-module such that M/qM has finite length. Then the integer \(d_{\mathfrak{q}}(\mathbf{M})\) is the dimension of the A-module M (§ 1, no. 4, def. 8).

We may assume \(\mathbf{M} \neq 0\). Let us first prove the inequality \(d_{\mathfrak{q}}(\mathbf{M}) \leqslant \dim_{\mathbf{A}}(\mathbf{M})\). By cor. 2 to prop. 9 of § 3, no. 5, there exists \(x_1, ..., x_r \in \mathfrak{q}\), with \(r = \dim_{\mathbf{A}}(\mathbf{M})\) and \(\text{long}(\mathbf{M}/\sum_{i=1}^{r} x_i \mathbf{M}) < +\infty\); set \(\mathfrak{x} = \sum_{i=1}^{r} x_i \mathbf{A}\). By prop. 4 of no. 3, we have \(d_{\mathfrak{x}}(\mathbf{M}) \leqslant r\); one has \(\mathfrak{x} \subset \mathfrak{q}\), whence \(\mathrm{H}_{\mathbf{M},\mathfrak{q}}^{(1)} \leqslant \mathrm{H}_{\mathbf{M},\mathfrak{x}}^{(1)}\) and therefore (lemma 2 of no. 1)

\[
d _ {\mathfrak{q}} (\mathbf {M}) \leqslant d _ {\mathfrak{x}} (\mathbf {M}) \leqslant r = \dim_ {\mathrm{A}} (\mathbf {M}).
\]

We now prove, by induction on \(\dim_{\mathbf{A}}(\mathbf{M})\) , the inequality \(\dim_{\mathbf{A}}(\mathbf{M}) \leqslant d_{\mathfrak{q}}(\mathbf{M})\) , obvious when \(\dim_{\mathbf{A}}(\mathbf{M}) = 0\) .

Suppose we have \(\dim_{\mathbf{A}}(\mathbf{M}) > 0\) , and \(\dim_{\mathbf{A}}(\mathbf{N}) \leqslant d_{\mathfrak{q}}(\mathbf{N})\) for every finitely generated A-module N such that \(\dim_{\mathbf{A}}(\mathbf{N}) < \dim_{\mathbf{A}}(\mathbf{M})\) . If \(0 = \mathbf{M}_0 \subset \mathbf{M}_1 \subset \ldots \subset \mathbf{M}_n = \mathbf{M}\) is a composition series of M, we have \(\dim_{\mathbf{A}}(\mathbf{M}) = \sup(\dim_{\mathbf{A}}(\mathbf{M}_i / \mathbf{M}_{i-1}))\) (§ 1, no. 4, prop. 9) and \(d_{\mathfrak{q}}(\mathbf{M}) = \sup(d_{\mathfrak{q}}(\mathbf{M}_i / \mathbf{M}_{i-1}))\) (No. 3, prop. 5). By IV, § 1, no. 4, th. 1, we may therefore assume that M is of the form A/p, where p is a prime ideal of A, and we have \(\mathfrak{p} \ne \mathfrak{m}_{A}\) because \(\dim_{\mathbf{A}}(\mathbf{M}) > 0\) . Let \(x \in \mathfrak{m}_{\mathbf{A}} - \mathfrak{p}\) ; the homothety \(x_{\mathbf{M}}\) of M = A/p is injective, and we have the exact sequence

\[
0 \longrightarrow \mathrm{M} \xrightarrow {x _ {\mathrm{M}}} \mathrm{M} \longrightarrow \mathrm{M} / x \mathrm{M} \longrightarrow 0.
\]

By § 3, no. 2, prop. 3, we have \(\dim_{\mathbf{A}}(\mathbf{M}/x\mathbf{M}) = \dim_{\mathbf{A}}(\mathbf{M}) - 1\); by prop. 5 of no. 3, and the preceding exact sequence, we have \(d_{\mathfrak{q}}(\mathbf{M}/x\mathbf{M}) \leqslant d_{\mathfrak{q}}(\mathbf{M}) - 1\). By the induction hypothesis, we therefore have

\[
\dim_ {\mathbf {A}} (\mathbf {M}) = \dim_ {\mathbf {A}} (\mathbf {M} / x \mathbf {M}) + 1 \leqslant d _ {\mathfrak {q}} (\mathbf {M} / x \mathbf {M}) + 1 \leqslant d _ {\mathfrak {q}} (\mathbf {M}),
\]

which completes the proof.

COROLLARY. --- Let A be a Noetherian ring, M a finitely generated A-module, and q an ideal of A such that M/qM has finite length. Then \(d_{\mathfrak{q}}(\mathbf{M})\) is the supremum of the dimensions \(\dim_{\mathrm{A}_{\mathfrak{m}}}\left(\mathrm{M}_{\mathfrak{m}}\right)\), where \(\mathfrak{m}\) ranges over the finite set \(S = \operatorname{Supp}(\mathbf{M}) \cap \mathrm{V}(\mathfrak{q})\), and \(e_{\mathfrak{q}}(\mathbf{M})\) is the sum of \(e_{\mathfrak{q}_{\mathfrak{m}}}\left(\mathrm{M}_{\mathfrak{m}}\right)\) over those \(\mathfrak{m}\in S\) for which \(\dim_{\mathrm{A}_{\mathfrak{m}}}\left(\mathrm{M}_{\mathfrak{m}}\right) = d_{\mathfrak{q}}(\mathbf{M})\).

For each integer \(n\), the length of \(\mathbf{M} / \mathfrak{q}^n\mathbf{M}\) is the sum of \(\mathrm{long}_{\mathbf{A}_{\mathfrak{m}}}(\mathbf{M}_{\mathfrak{m}} / \mathfrak{q}_{\mathfrak{m}}^n\mathbf{M}_{\mathfrak{m}})\) (IV, § 2, no. 5, cor. 1 to prop. 7 and corollary to prop. 8). Consequently, we have \(\mathbf{H}_{\mathbf{M},\mathfrak{q}} = \sum_{\mathfrak{m}\in S}\mathbf{H}_{\mathbf{M}_{\mathfrak{m}},\mathfrak{q}_{\mathfrak{m}}}\), whence the corollary.

Remarks. --- 1) We also have \(d_{\mathfrak{q}}(\mathbf{M}) = \sup_{\mathfrak{m} \in \mathbf{V}(\mathfrak{q})} \dim(\mathbf{M}_{\mathfrak{m}})\) , that is, \(d_{\mathfrak{q}}(\mathbf{M}) = \dim(\hat{\mathbf{M}})\) , where \(\hat{\mathbf{M}}\) is the completion of \(\mathbf{M}\) for the q-adic topology (§ 3, No. 4, prop. 8).

\begin{enumerate}
\def\labelenumi{\arabic{enumi})}
\setcounter{enumi}{1}
\tightlist
\item
  Suppose q is contained in the radical of A; then \(\dim(\tilde{\mathbf{M}})=\dim(\mathbf{M})\) (loc. cit., cor. 1), hence \(d_{\mathfrak{q}}(\mathbf{M})=\dim(\mathbf{M})\) .
\end{enumerate}

\subsection{Hilbert-Samuel series of a quotient module}

Lemma 5. — Let A be a ring, M an A-module, and \((\mathbf{P}_{n})\) , \((\mathbf{Q}_{n})\) two descending filtrations on M consisting of submodules. Suppose that we have \(P_{n} \supset Q_{n}\) and \(\text{long}_{\mathbf{A}}(\mathbf{P}_{n}/\mathbf{Q}_{n}) < +\infty\) for every \(n \in Z\) and that there exists an integer \(n_{1}\) such that \(Q_{n_{1}} = M\). In \(\mathbf{Z}((T))\), we have the inequalities

\[
\begin{array}{r l} \sum_ {n \in \mathbf {Z}} \mathrm{long} _ {\mathbf {A}} ((\mathrm{P} _ {n + 1} \cap \mathrm{Q} _ {n}) / \mathrm{Q} _ {n + 1}). & \mathrm{T} ^ {n} \leqslant \sum_ {n \in \mathbf {Z}} \mathrm{long} _ {\mathbf {A}} (\mathrm{P} _ {n + 1} / \mathrm{Q} _ {n + 1}). \\ & \mathrm{T} ^ {n} \leqslant \\ & \leqslant (1 - \mathrm{T}) ^ {- 1} \sum_ {n \in \mathbf {Z}} \mathrm{long} _ {\mathbf {A}} ((\mathrm{P} _ {n + 1} \cap \mathrm{Q} _ {n}) / \mathrm{Q} _ {n + 1}). \end{array}
\]

We need to prove that we have the inequalities

\begin{enumerate}
\def\labelenumi{(\arabic{enumi})}
\setcounter{enumi}{14}
\tightlist
\item
\end{enumerate}

\[
\operatorname{long} \left(\left(\mathrm{P} _ {n + 1} \cap \mathrm{Q} _ {n}\right) / \mathrm{Q} _ {n + 1}\right) \leqslant \operatorname{long} \left(\mathrm{P} _ {n + 1} / \mathrm{Q} _ {n + 1}\right),\tag{16}
\]

\[
\operatorname{long} \left(\mathrm{P} _ {n + 1} / \mathrm{Q} _ {n + 1}\right) \leqslant \sum_ {i \leqslant n} \operatorname{long} \left(\left(\mathrm{P} _ {i + 1} \cap \mathrm{Q} _ {i}\right) / \mathrm{Q} _ {i + 1}\right).
\]

The first is evident. On the other hand, we have \(P_{n+1} \cap Q_i = P_{n+1}\) for \(i \leqslant n_1\) and \(P_{n+1} \cap Q_{n+1} = Q_{n+1}\) ; we deduce the inequality

\[
\operatorname{long} \left(\mathrm{P} _ {n + 1} / \mathrm{Q} _ {n + 1}\right) \leqslant \sum_ {i \leqslant n} \operatorname{long} \left(\left(\mathrm{P} _ {n + 1} \cap \mathrm{Q} _ {i}\right) / \left(\mathrm{P} _ {n + 1} \cap \mathrm{Q} _ {i + 1}\right)\right).
\]

But the A-module \((\mathrm{P}_{n+1} \cap \mathrm{Q}_i)/(\mathrm{P}_{n+1} \cap \mathrm{Q}_{i+1})\) is isomorphic to a submodule of \((\mathrm{P}_{i+1} \cap \mathrm{Q}_i)/\mathrm{Q}_{i+1}\) , and inequality (16) follows.

Lemma 6. --- Let A be a ring, M an A-module, and \((F_{n})\) a decreasing filtration on M consisting of submodules; suppose there exists an integer \(n_{1}\) such that \(F_{n_{1}} = M\). Let f be an endomorphism of M, \(M'\) its kernel, and \(M''\) its cokernel. We endow \(M'\) with the filtration \((F_{n}')\) induced by \((F_{n})\) and \(M''\) with the quotient filtration \((F_{n}'')\) of \((F_{n})\). We assume that \(F_{n}/F_{n+1}\) is of finite length for every \(n \in \mathbf{Z}\) and that there exists an integer \(\delta\) such that \(f(F_{n}) \subset F_{n+\delta}\). Let \(\varphi\) be the graded endomorphism of degree \(\delta\) of the graded module \(\operatorname{gr}(M) = \bigoplus_{n} F_{n}/F_{n+1}\) induced by f. Between the following elements of \(\mathbf{Z}((T))\)

\[
\mathrm{H} _ {\mathbf {M}} = \sum_ {n \in \mathbf {Z}} \operatorname{long} _ {\mathbf {A}} (\mathrm{F} _ {n} / \mathrm{F} _ {n + 1}). \mathrm{T} ^ {n}
\]

\[
\mathrm{H} _ {\mathbf {M} ^ {\prime}} = \sum_ {n \in \mathbf {Z}} \operatorname{long} _ {\mathbf {A}} (\mathrm{F} _ {n} ^ {\prime} / \mathrm{F} _ {n + 1} ^ {\prime}). \mathrm{T} ^ {n}
\]

\[
\mathrm{H} _ {\mathrm{M} ^ {\prime \prime}} = \sum_ {n \in \mathbf {Z}} \operatorname{long} _ {\mathrm{A}} (\mathrm{F} _ {n} ^ {\prime \prime} / \mathrm{F} _ {n + 1} ^ {\prime \prime}). \mathrm{T} ^ {n}
\]

\[
\mathrm{P} _ {\operatorname{Ker} (\varphi)} = \sum_ {n \in \mathbf {Z}} \operatorname{long} _ {\mathrm{A}} (\operatorname{Ker} (\varphi_ {n})). \mathrm{T} ^ {n},
\]

we have the inequalities

\begin{enumerate}
\def\labelenumi{(\arabic{enumi})}
\setcounter{enumi}{16}
\tightlist
\item
\end{enumerate}

\[
\mathrm{H} _ {\mathbf {M} ^ {\prime}} \leqslant \mathrm{P} _ {\mathbf {K e r} (\varphi)}\tag{18}
\]

\[
(1 - \mathrm{T} ^ {\delta}). \mathrm{H} _ {\mathbf {M}} ^ {(1)} + \mathrm{T} ^ {\delta}. \mathrm{P} _ {\operatorname{Ker} (\varphi)} \leqslant \mathrm{H} _ {\mathbf {M} ^ {\prime \prime}} ^ {(1)} \leqslant (1 - \mathrm{T} ^ {\delta}). \mathrm{H} _ {\mathbf {M}} ^ {(1)} + \mathrm{T} ^ {\delta}. \mathrm{P} _ {\operatorname{Ker} (\varphi)} ^ {(1)}.
\]

The sequence of submodules \(\mathbf{G}_{n}=f^{-1}(\mathbf{F}_{n+\delta})\) on M is a decreasing filtration, and we have \(F_{n}\subset G_{n}\) for every integer n.

By definition, we have \(\operatorname{Ker}(\varphi_n) = (\mathbf{G}_{n+1} \cap \mathbf{F}_n)/\mathbf{F}_{n+1}\) , whence

\[
\mathrm{P} _ {\operatorname{Ker} (\varphi)} = \sum_ {n \in \mathbf {Z}} \operatorname{long} _ {\mathrm{A}} ((\mathrm{G} _ {n + 1} \cap \mathrm{F} _ {n}) / \mathrm{F} _ {n + 1}). \mathrm{T} ^ {n}.\tag{19}
\]

For every \(n\) , the A-module \((\mathbf{M}' \cap \mathbf{F}_n) / (\mathbf{M}' \cap \mathbf{F}_{n+1})\) is identified with a submodule of \((\mathbf{G}_{n+1} \cap \mathbf{F}_n) / \mathbf{F}_{n+1}\) , and inequality (17) follows immediately from (19). Moreover, by Lemma 5, we have

\[
\mathrm{P} _ {\operatorname{Ker} (\varphi)} \leqslant \sum_ {n \in \mathbf {Z}} \operatorname{long} _ {\mathrm{A}} \left(\mathrm{G} _ {n + 1} / \mathrm{F} _ {n + 1}\right). \mathrm{T} ^ {n} \leqslant \mathrm{P} _ {\operatorname{Ker} (\varphi)} ^ {(1)}.\tag{20}
\]

For every \(n \in \mathbf{Z}\) we have an exact sequence of A-modules

\[
0 \longrightarrow \mathrm{G} _ {n + 1} / \mathrm{F} _ {n + 1} \longrightarrow \mathrm{M} / \mathrm{F} _ {n + 1} \xrightarrow {f _ {n}} \mathrm{M} / \mathrm{F} _ {n + \delta + 1} \longrightarrow \mathrm{M} ^ {\prime \prime} / \mathrm{F} _ {n + \delta + 1} ^ {\prime \prime} \longrightarrow 0,
\]

where \(f_{n}\) is deduced from \(f\) by passing to quotients. Consequently, we have

\[
\operatorname{long} _ {\mathrm{A}} \left(\mathrm{M} ^ {\prime \prime} / \mathrm{F} _ {n + \delta + 1} ^ {\prime \prime}\right) = \operatorname{long} _ {\mathrm{A}} \left(\mathrm{M} / \mathrm{F} _ {n + \delta + 1}\right) - \operatorname{long} _ {\mathrm{A}} \left(\mathrm{M} / \mathrm{F} _ {n + 1}\right) + \operatorname{long} _ {\mathrm{A}} \left(\mathrm{G} _ {n + 1} / \mathrm{F} _ {n + 1}\right).
\]

Multiplying by \(T^{n+\delta}\) and summing over n, we obtain

\[
\mathrm{H} _ {\mathbf {M} ^ {\prime \prime}} ^ {(1)} = (1 - \mathrm{T} ^ {\delta}) \mathrm{H} _ {\mathbf {M}} ^ {(1)} + \mathrm{T} ^ {\delta}. \sum_ {n \in \mathbf {Z}} \operatorname{long} _ {\mathrm{A}} (\mathrm{G} _ {n + 1} / \mathrm{F} _ {n + 1}). \mathrm{T} ^ {n},\tag{21}
\]

and inequality (18) follows immediately from (20) and (21).

Lemma 7. — We keep the notations of Lemma 6.

\begin{enumerate}
\def\labelenumi{\alph{enumi})}
\item
  We have the inequality \(\mathbf{H}_{\mathbf{M}^{\prime \prime}}^{(1)}\geqslant \frac{1 - \mathrm{T}^{\delta}}{1 - \mathrm{T}}\mathbf{H}_{\mathbf{M}}.\)
\item
  For equality to hold, it is necessary and sufficient that φ be injective.
\item
  If this is the case, we have \(\mathbf{M}^{\prime}\subset \bigcap_{n}\mathbf{F}_{n}\) and the sequence of A-modules
\end{enumerate}

\[
0 \longrightarrow \operatorname{gr} (\mathrm{M}) \stackrel {\varphi} {\longrightarrow} \operatorname{gr} (\mathrm{M}) \stackrel {v} {\longrightarrow} \operatorname{gr} (\mathrm{M} ^ {\prime \prime}) \longrightarrow 0,
\]

where v is the canonical map, is exact.

Assertions a) and b) follow from formula (18) of Lemma 6 and from the definition \(\mathbf{H}_{\mathbf{M}}^{(1)} = (1 - \mathbf{T})^{-1}.\mathbf{H}_{\mathbf{M}}\) .

Suppose that \(\varphi\) is injective. By III, § 2, no. 8, Th. 1, (i), we have

\[
\operatorname{Ker} (f) \subset f ^ {- 1} (\mathrm{F} _ {n + \delta}) = \mathrm{F} _ {n}
\]

for all n, whence the first assertion of c). Moreover, we have an exact sequence

\[
0 \longrightarrow \mathrm{M} / \mathrm{M} ^ {\prime} \stackrel {f ^ {\prime}} {\longrightarrow} \mathrm{M} \longrightarrow \mathrm{M} / f (\mathrm{M}) \longrightarrow 0,
\]

where \(f'\) is deduced from \(f\) by passing to the quotient. If \(\varphi\) is injective, we have as above \(f'^{-1}(\mathbf{F}_n) = \mathbf{F}_{n-\delta}/\mathbf{M}'\). Thus the filtration on \(\mathbf{M}/\mathbf{M}'\) deduced by inverse image under \(f'\) of the filtration \(\mathbf{F}\) on \(\mathbf{M}\) is the filtration \(n \mapsto \mathbf{F}_{n-\delta}/\mathbf{M}'\); the associated graded is \(\operatorname{gr}(\mathbf{M})(-\delta)\) and we have an exact sequence of graded modules (III, § 2, no. 4, prop. 2)

\[
0 \longrightarrow \operatorname{gr} (\mathrm{M}) (- \delta) \xrightarrow {\varphi^ {\prime}} \operatorname{gr} (\mathrm{M}) \longrightarrow \operatorname{gr} \left(\mathrm{M} ^ {\prime \prime}\right) \longrightarrow 0,
\]

where \(\varphi_{n}^{\prime} = \varphi_{n - \delta}\) for every \(n\). This completes the proof of \(c\) ).

PROPOSITION 6. — Let A be a Noetherian ring, M a finitely generated A-module, and q an ideal of A such that M/qM has finite length. Let F be a q-good filtration of M, and let gr(A) = \(\bigoplus_{n\geqslant 0}(\mathfrak{q}^n /\mathfrak{q}^{n + 1})\) the graded ring associated with A for the q-adic filtration.

Let \((x_{1},\ldots,x_{s})\) be a sequence of elements of A, \((\delta_{1},\ldots,\delta_{s})\) a sequence of strictly positive integers such that \(x_{i}\in\mathfrak{q}^{\delta_{i}}\) for \(1\leqslant i\leqslant s\) , and let \(\xi_{i}\) be the class of \(x_{i}\) in \(\mathrm{gr}_{\delta_{i}}(\mathbf{A})=\mathfrak{q}^{\delta_{i}}/\mathfrak{q}^{\delta_{i}+1}\) .

\begin{enumerate}
\def\labelenumi{\alph{enumi})}
\tightlist
\item
  Endow the A-module \(\overline{\mathbf{M}} = \mathbf{M}/(x_1\mathbf{M} + \cdots + x_s\mathbf{M})\) with the quotient q-good filtration \(\overline{\mathbf{F}}\) of F. We then have in \(\mathbf{Z}((\mathbf{T}))\) the inequality
\end{enumerate}

\[
\mathrm{H} _ {\mathrm{M}, \overline {{\mathrm{F}}}} ^ {(s)} \geqslant \left(\prod_ {i = 1} ^ {s} \frac {1 - \mathrm{T} ^ {\delta_ {i}}}{1 - \mathrm{T}}\right). \mathrm{H} _ {\mathrm{M}, \mathrm{F}}.\tag{22}
\]

\begin{enumerate}
\def\labelenumi{\alph{enumi})}
\setcounter{enumi}{1}
\tightlist
\item
  For equality to hold in (22), it is necessary and sufficient that the sequence \((\xi_1, \ldots, \xi_s)\) of elements of the ring gr(A) be completely secant for the module gr(M) = \(\bigoplus_{n} (\mathbf{F}_n / \mathbf{F}_{n+1})\) .
\end{enumerate}

In this case, the canonical homomorphism from \(\mathrm{gr}(\mathbf{M}) / \sum_{i=1}^{s} \xi_i \cdot \mathrm{gr}(\mathbf{M})\) into \(\mathrm{gr}(\overline{\mathbf{M}}) = \bigoplus_{n} (\overline{\mathbf{F}}_n / \overline{\mathbf{F}}_{n+1})\) is an isomorphism.

\begin{enumerate}
\def\labelenumi{\alph{enumi})}
\setcounter{enumi}{2}
\tightlist
\item
  Suppose the conditions of b) are satisfied, and that each of the A-modules \(\mathbf{M}_{i} = \mathbf{M}/(x_{1}\mathbf{M} + \cdots + x_{i}\mathbf{M})\) for \(0 \leqslant i < s\) is separated for the q-adic topology \footnote{This occurs in particular if \(q\) is contained in the radical of \(A\) (III, § 3, no. 3, Prop. 6).}. Then the sequence \((x_{1}, ..., x_{s})\) is completely secant for the A-module M.
\end{enumerate}

When \(s = 1\) , one has \(\bigcap_{n} F_{n} = \bigcap_{n} q^{n} M\) and the sequence \(\{\xi_{1}\}\) is completely secant for \(\operatorname{gr}(M)\) if and only if the homothety with ratio \(\xi_{1}\) in \(\operatorname{gr}(M)\) is injective. Prop. 6 then follows immediately from Lemma 7 applied to the homothety \(f = (x_{1})_{M}\) in \(M\) .

Suppose that we have \(s \geqslant 2\) and reason by induction on \(s\). The induction hypothesis applied to the A-module \(\mathbf{M}_{1} = \mathbf{M}/x_{1}\mathbf{M}\) endowed with the quotient filtration G of F, and to the sequence \((x_{2}, ..., x_{s})\) yields the inequality

\[
\mathrm{H} _ {\overline {{\mathrm{M,F}}}} ^ {(s - 1)} \geqslant \left(\prod_ {i = 2} ^ {s} \frac {1 - \mathrm{T} ^ {\delta_ {i}}}{1 - \mathrm{T}}\right). \mathrm{H} _ {\mathrm{M} _ {1}, \mathrm{G}};\tag{23}
\]

there is equality if and only if the sequence \((\xi_{2}, \ldots, \xi_{s})\) is completely secant for the \(\operatorname{gr}(A)\)-module \(\operatorname{gr}(\mathbf{M}_{1}) = \bigoplus_{n} \mathrm{G}_{n}/\mathrm{G}_{n+1}\). Since the elements \(\frac{1 - \mathrm{T}^{\delta_{i}}}{1 - \mathrm{T}}\) of \(\mathbf{Z}((\mathrm{T}))\) are positive, the case \(s = 1\) already treated and formula (23) yield the inequalities

\[
\mathrm{H} _ {\mathrm{M}, \mathrm{F}} ^ {(s)} \geqslant \left(\prod_ {i = 2} ^ {s} \frac {1 - \mathrm{T} ^ {\delta_ {i}}}{1 - \mathrm{T}}\right). \mathrm{H} _ {\mathrm{M} _ {1}, \mathrm{G}} ^ {(1)} \geqslant \left(\prod_ {i = 1} ^ {s} \frac {1 - \mathrm{T} ^ {\delta_ {i}}}{1 - \mathrm{T}}\right). \mathrm{H} _ {\mathrm{M}, \mathrm{F}}.\tag{24}
\]

This proves a).

One can have equality in (22) only if one has simultaneously equality in (23) and equality

\[
\mathrm{H} _ {\mathbf {M} _ {1}, \mathbf {G}} ^ {(1)} = \left(\frac {1 - \mathrm{T} ^ {\delta_ {1}}}{1 - \mathrm{T}}\right) \cdot \mathrm{H} _ {\mathbf {M}, \mathbf {F}}.\tag{25}
\]

This latter relation means that \(\{\xi_{1}\}\) is completely secant for \(\operatorname{gr}(M)\) and implies that the canonical homomorphism from \(\operatorname{gr}(M)/\xi_{1}\operatorname{gr}(M)\) into \(\operatorname{gr}(M_{1})\) is an isomorphism. In other words, equality holds in (22) if and only if \(\{\xi_{1}\}\) is completely secant for \(\operatorname{gr}(M)\) and \(\{\xi_{2}, ..., \xi_{s}\}\) is completely secant for \(\operatorname{gr}(M)/\xi_{1}\operatorname{gr}(M)\). This means that \(\{\xi_{1}, ..., \xi_{s}\}\) is completely secant for \(\operatorname{gr}(M)\) (A, X, p. 160, cor. 2). We have thus proved the equivalence of the two conditions of \(b\)). Assume they are satisfied; then, by the induction hypothesis, \(\operatorname{gr}(M)\) is identified with \(\operatorname{gr}(M_{1})/\sum_{i=2}^{s} \xi_{i}\operatorname{gr}(M_{1})\); moreover, since \(\operatorname{gr}(M_{1})\) is identified with \(\operatorname{gr}(M)/\xi_{1}\operatorname{gr}(M)\), the last assertion of \(b\)) is thus satisfied.

Suppose now that \(\{\xi_1, \ldots, \xi_s\}\) is completely secant for \(\operatorname{gr}(\mathbf{M})\) and that each \(\mathbf{M}_i\) is separated for the q-adic topology (for \(0 \leqslant i < s\) ). From the above and the induction hypothesis, the sequence \((x_2, \ldots, x_s)\) is completely secant for \(\mathbf{M}_1\); since we have \(\mathbf{M}_1 = \mathbf{M}/x_1\mathbf{M}\) and \(\{x_1\}\) is completely secant for \(\mathbf{M}\), the sequence \((x_1, x_2, \ldots, x_s)\) is completely secant for \(\mathbf{M}\) (A, X, p. 160, th. 1).

\section{REGULAR LOCAL RINGS}
